\documentclass[12pt,a4paper]{article}
\usepackage[utf8]{inputenc}
\usepackage[T1]{fontenc}
\usepackage[ngerman]{babel}
\usepackage{geometry}
\usepackage{graphicx}
\usepackage{listings}
\usepackage{xcolor}
\usepackage{tikz}
\usepackage{booktabs}
\usepackage{array}
\usepackage{tabularx}
\usepackage{amsmath}
\usepackage{hyperref}
\usepackage{fancyhdr}
\usepackage{parskip}
\usepackage{enumitem}
\usepackage{tcolorbox}
\usepackage{float}
\usepackage{mdframed}

\usetikzlibrary{shapes.geometric, arrows.meta, positioning, fit, backgrounds, calc}

\geometry{margin=2.5cm}

\hypersetup{
    colorlinks=true,
    linkcolor=blue!60!black,
    urlcolor=blue!60!black,
}

\definecolor{codebg}{RGB}{245,245,245}
\definecolor{codeframe}{RGB}{200,200,200}
\definecolor{javakw}{RGB}{0,0,200}
\definecolor{javastr}{RGB}{0,128,0}
\definecolor{javacomment}{RGB}{128,128,128}
\definecolor{mermbg}{RGB}{240,248,255}
\definecolor{mermframe}{RGB}{70,130,180}

\lstset{
    language=Java,
    backgroundcolor=\color{codebg},
    frame=single,
    rulecolor=\color{codeframe},
    basicstyle=\ttfamily\small,
    keywordstyle=\color{javakw}\bfseries,
    stringstyle=\color{javastr},
    commentstyle=\color{javacomment}\itshape,
    breaklines=true,
    showstringspaces=false,
    tabsize=4,
    numbers=left,
    numberstyle=\tiny\color{gray},
    numbersep=8pt,
}

\lstdefinestyle{mermaid}{
    language={},
    backgroundcolor=\color{mermbg},
    frame=single,
    rulecolor=\color{mermframe},
    basicstyle=\ttfamily\small,
    breaklines=true,
    showstringspaces=false,
    numbers=none,
    morekeywords={classDiagram,sequenceDiagram,flowchart,class,participant,
                  loop,end,alt,else,opt,Note,TD,LR,TB,
                  direction,subgraph,style,linkStyle},
    keywordstyle=\color{blue!70!black}\bfseries,
}

\pagestyle{fancy}
\fancyhf{}
\rhead{Eisenbahn-Fahrplan Simulation}
\lhead{IHK Abschlussprojekt}
\rfoot{Seite \thepage}

\tcbuselibrary{skins, breakable}

\newtcolorbox{beispielbox}[1][]{
    colback=blue!5,
    colframe=blue!40!black,
    fonttitle=\bfseries,
    title=Beispiel,
    #1
}

\newtcolorbox{hinweisbox}[1][]{
    colback=orange!5,
    colframe=orange!60!black,
    fonttitle=\bfseries,
    title=Hinweis,
    #1
}

\newtcolorbox{designbox}[1][]{
    colback=green!5,
    colframe=green!50!black,
    fonttitle=\bfseries,
    title=Designentscheidung,
    #1
}

% ─────────────────────────────────────────
\begin{document}

\begin{titlepage}
    \centering
    \vspace*{2cm}
    {\Huge\bfseries Eisenbahn-Fahrplan Simulation\par}
    \vspace{0.5cm}
    {\Large Technische Dokumentation\par}
    \vspace{2cm}
    \rule{\linewidth}{0.4pt}\\[0.5cm]
    {\large IHK Abschlussprojekt -- MATSE\par}
    \rule{\linewidth}{0.4pt}
    \vspace{2cm}
    {\large\itshape Kollisionserkennung, Wartezeitermittlung\\und Umsetzung der drei Fahrplanstrategien\par}
    \vfill
    {\large 2026\par}
\end{titlepage}

\tableofcontents
\newpage

% ─────────────────────────────────────────
\section{Aufgabenbeschreibung}

Das Programm simuliert den Fahrplan einer eingleisigen Eisenbahnstrecke zwischen mehreren Bahnhöfen. Auf dieser Strecke fährt ein Zug von Bahnhof \textbf{A} nach Bahnhof \textbf{Z} (Hinfahrt) und anschließend zurück (Rückfahrt). Da die Strecke nur \textbf{eingleisig} ist, darf jeweils nur ein Zug gleichzeitig auf einem Streckenabschnitt zwischen zwei aufeinanderfolgenden Bahnhöfen fahren. Andernfalls entsteht eine \textbf{Kollision}.

Ziel ist es, Fahrpläne nach drei verschiedenen Strategien zu berechnen:
\begin{enumerate}
    \item \textbf{Einfache Fahrt}: Hin- und Rückfahrt werden ohne Rücksicht auf Kollisionen berechnet.
    \item \textbf{Einseitiges Warten}: Nur der Rückzug wartet an Kollisionspunkten.
    \item \textbf{Beidseitiges Warten}: Die Wartezeit wird auf beide Züge verteilt, um die Strafe zu minimieren.
\end{enumerate}

Die \textbf{Strafe} für eine Strategie wird als Summe der quadrierten Wartezeiten berechnet, um lange Wartezeiten stärker zu bestrafen als viele kurze.

Die Eingabe erfolgt über Textdateien, die Strecke, Abstände zwischen den Bahnhöfen sowie die Startzeit der Hinfahrt enthalten. Die Ausgabe wird sowohl auf der Konsole angezeigt als auch in Dateien gespeichert.

% ─────────────────────────────────────────
\section{Modellierung der Datenstrukturen}

\subsection{Überblick}

Das Programm besteht aus folgenden Hauptklassen:

\begin{center}
\begin{tabular}{lll}
\toprule
\textbf{Klasse / Record} & \textbf{Paket} & \textbf{Aufgabe} \\
\midrule
\texttt{BahnhofNode} & \texttt{model} & Einzelner Bahnhof mit Zeitdaten \\
\texttt{BahnhofPlan} & \texttt{model} & Verkettete Liste aller Bahnhöfe \\
\texttt{SimulationsDaten} & \texttt{model} & Eingabedaten (Plan + Startzeit) \\
\texttt{FahrplanStrategie} & \texttt{strategie} & Abstrakte Basisklasse (Strategy Pattern) \\
\texttt{EinfacheFahrtStrategie} & \texttt{strategie} & Konkrete Strategie 1 \\
\texttt{EinseitigesWartenStrategie} & \texttt{strategie} & Konkrete Strategie 2 \\
\texttt{BeidseitigesWartenStrategie} & \texttt{strategie} & Konkrete Strategie 3 \\
\texttt{EingabeLeser} & \texttt{io} & Einlesen der Textdatei \\
\texttt{AusgabeManager} & \texttt{io} & Formatierte Ausgabe \\
\bottomrule
\end{tabular}
\end{center}

\subsection{BahnhofNode -- Der Knoten der verketteten Liste}

Jeder Bahnhof wird durch ein Objekt der Klasse \texttt{BahnhofNode} repräsentiert. Der Knoten speichert sowohl seine eigenen Zeitdaten als auch Zeiger auf den Vorgänger- und Nachfolgeknoten:

\begin{lstlisting}[caption={Felder von BahnhofNode}]
private String name;
private Integer ankunfthin;    // Ankunftszeit Hinfahrt
private Integer abfahrthin;    // Abfahrtszeit Hinfahrt
private Integer wartezeitHin = 0;
private Integer ankunftrueck;  // Ankunftszeit Rueckfahrt
private Integer abfahrtrueck;  // Abfahrtszeit Rueckfahrt
private Integer wartezeitRueck = 0;
private BahnhofNode next;      // Zeiger auf naechsten Bahnhof
private BahnhofNode prev;      // Zeiger auf vorherigen Bahnhof
private int distanceToNext;    // Fahrtzeit zum naechsten Bahnhof (Minuten)
\end{lstlisting}

\subsection{BahnhofPlan -- Die doppelt verkettete Liste}

Die Klasse \texttt{BahnhofPlan} verwaltet alle Bahnhöfe als \textbf{doppelt verkettete Liste}. Sie hält direkte Zeiger auf den ersten Knoten (\texttt{head}) und den letzten Knoten (\texttt{tail}):

\begin{figure}[H]
\centering
\begin{tikzpicture}[
    node distance=0mm,
    box/.style={draw, minimum width=2.2cm, minimum height=0.7cm, font=\small\ttfamily},
    ptr/.style={-Stealth, thick},
]
\node[box, fill=blue!10] (A) {A};
\node[box, fill=blue!10, right=1.5cm of A] (B) {B};
\node[box, fill=blue!10, right=1.5cm of B] (C) {C};
\node[box, fill=blue!10, right=1.5cm of C] (D) {\ldots};

\node[above=0.4cm of A, font=\small\bfseries] (headlbl) {head};
\draw[ptr] (headlbl) -- (A);
\node[above=0.4cm of C, font=\small\bfseries] (taillbl) {tail};
\draw[ptr] (taillbl) -- (C);

\draw[ptr, red!70!black] (A.east) -- node[above, font=\tiny] {next} (B.west);
\draw[ptr, red!70!black] (B.east) -- node[above, font=\tiny] {next} (C.west);
\draw[ptr, red!70!black] (C.east) -- node[above, font=\tiny] {next} (D.west);

\draw[ptr, green!50!black, bend right=30] (B.south) to node[below, font=\tiny] {prev} (A.south);
\draw[ptr, green!50!black, bend right=30] (C.south) to node[below, font=\tiny] {prev} (B.south);
\draw[ptr, green!50!black, bend right=30] (D.south) to node[below, font=\tiny] {prev} (C.south);

\node[below=0.9cm of A, xshift=0.75cm, font=\tiny, red!70!black] {dist=5};
\node[below=0.9cm of B, xshift=0.75cm, font=\tiny, red!70!black] {dist=8};
\node[below=0.9cm of C, xshift=0.75cm, font=\tiny, red!70!black] {dist=?};
\end{tikzpicture}
\caption{Doppelt verkettete Liste der Bahnhöfe (Beispiel: A--B--C--\ldots)}
\end{figure}

\subsubsection{Warum eine doppelt verkettete Liste?}

Die doppelt verkettete Liste ist für dieses Problem besonders geeignet, weil:

\begin{itemize}
    \item \textbf{Hinfahrt}: Der Algorithmus läuft von \texttt{head} nach \texttt{tail} und nutzt dabei \texttt{getNext()} -- eine natürliche Vorwärtstraversierung.
    \item \textbf{Rückfahrt}: Der Algorithmus läuft von \texttt{tail} nach \texttt{head} und nutzt \texttt{getPrev()} -- eine natürliche Rückwärtstraversierung.
    \item \textbf{Lokale Kollisionsprüfung}: \texttt{isKollision()} benötigt nur den aktuellen Knoten und \texttt{this.next} -- es ist kein globaler Zugriff auf den gesamten Plan notwendig.
    \item \textbf{Dynamische Länge}: Die Liste wächst beim Einlesen der Daten automatisch, ohne dass vorab die Anzahl der Bahnhöfe bekannt sein muss.
\end{itemize}

\subsection{SimulationsDaten -- Java Record}

\begin{lstlisting}[caption={SimulationsDaten als Record}]
public record SimulationsDaten(BahnhofPlan plan, int startZeit) {}
\end{lstlisting}

Ein Java-\texttt{record} ist ein kompakter unveränderlicher Datenträger. Er kapselt den gelesenen \texttt{BahnhofPlan} und die Startzeit der Hinfahrt in einem einzigen Objekt, das \texttt{EingabeLeser} zurückgibt und \texttt{Main} weiterverarbeitet.

\subsection{Das Strategy Pattern}

Die abstrakte Klasse \texttt{FahrplanStrategie} bildet die gemeinsame Schnittstelle aller drei Strategien:

\begin{lstlisting}[caption={Abstrakte Basisklasse FahrplanStrategie}]
public abstract class FahrplanStrategie {
    public abstract String getStrategieName();
    public abstract void berechneFahrplan(BahnhofPlan plan, int startZeit);
    public abstract int getSummeStrafen();
    public abstract int getSummeWarteZeit();
    protected BahnhofPlan aktuellerPlan;
    // Gemeinsame statische Hilfsmethoden fuer alle Strategien ...
}
\end{lstlisting}

Der Vorteil einer abstrakten Klasse gegenüber einem Interface liegt darin, dass gemeinsam genutzte Hilfsmethoden (\texttt{hinfahrtOhneKollision}, \texttt{rueckfahrtMitKollision}, \texttt{beidseitigesWarten} usw.) direkt dort als \texttt{static}-Methoden implementiert werden können. Alle konkreten Strategien erben diese Methoden und rufen sie auf.

In \texttt{Main.java} werden alle drei Strategien gleichartig behandelt:

\begin{lstlisting}[caption={Verwendung des Strategy Patterns in Main}]
FahrplanStrategie[] strategien = {
    new EinfacheFahrtStrategie(),
    new EinseitigesWartenStrategie(),
    new BeidseitigesWartenStrategie()
};
for (FahrplanStrategie strategie : strategien) {
    strategie.berechneFahrplan(plan, startZeit);
    gesamtausgabe.append(AusgabeManager.druckeErgebnis(strategie, plan));
}
\end{lstlisting}

Der aufrufende Code benötigt keine \texttt{if-else}-Kette. Neue Strategien können einfach hinzugefügt werden, indem man \texttt{FahrplanStrategie} erweitert und das neue Objekt ins Array aufnimmt.

\subsection{Geteilte Objektreferenz: Strategien bauen aufeinander auf}
\label{sec:referenz}

\begin{designbox}
Alle drei Strategien erhalten \textbf{dieselbe} \texttt{BahnhofPlan}-Referenz. Es wird \textbf{keine Kopie} des Objekts angelegt. Die Strategien schreiben ihre Ergebnisse direkt in dieselben \texttt{BahnhofNode}-Objekte -- und lesen die Ergebnisse der vorherigen Strategie von dort.
\end{designbox}

Dies ist eine bewusste Designentscheidung mit folgenden Konsequenzen:

\textbf{Strategie 1 -- EinfacheFahrtStrategie}:
Berechnet sowohl die \emph{Hinfahrt} als auch die \emph{Rückfahrt} und schreibt alle Zeitwerte (\texttt{ankunfthin}, \texttt{abfahrthin}, \texttt{ankunftrueck}, \texttt{abfahrtrueck}) in die Knoten des gemeinsamen Plans.

\textbf{Strategie 2 -- EinseitigesWartenStrategie}:
Liest die bereits von Strategie~1 gesetzten Hinfahrtdaten (\texttt{ankunfthin}, \texttt{abfahrthin}) aus den Knoten, ohne sie neu zu berechnen:

\begin{lstlisting}[caption={EinseitigesWartenStrategie baut auf Hinfahrtdaten von Strategie 1 auf}]
public void berechneFahrplan(BahnhofPlan plan, int startZeit) {
    this.aktuellerPlan = plan;
    // Liest ankunfthin(tail) -- wurde von EinfacheFahrtStrategie gesetzt!
    int startZeitRueck = (aktuellerPlan.getTail().getAnkunfthin() + 1) % 60;
    rueckfahrtMitKollision(aktuellerPlan, startZeitRueck);
}
\end{lstlisting}

Die Zeile \texttt{getTail().getAnkunfthin()} greift auf den Wert zu, den \texttt{EinfacheFahrtStrategie} zuvor in den letzten Knoten geschrieben hat. Die Hinfahrt wird \emph{nicht} neu berechnet.

\textbf{Strategie 3 -- BeidseitigesWartenStrategie}:
Funktioniert analog. Sie liest ebenfalls \texttt{ankunfthin} der Hinfahrt (gesetzt durch Strategie~1) und überschreibt im Laufe ihrer Berechnung die Rückfahrtdaten sowie ggf. auch die Hinwartezeiten.

\textbf{Warum kein Reset / keine Kopie?}

Die Hinfahrtdaten (gesetzt von Strategie~1) sind für alle drei Strategien identisch, weil Strategien 2 und 3 die Hinfahrt unverändert lassen -- sie verändern nur die Rückfahrtzeiten und Wartezeiten. Dadurch ist es korrekt, dieselbe Referenz weiterzugeben, ohne das Objekt vorher zurückzusetzen. Lediglich die Rückfahrtwerte und Wartezeiten werden bei jeder Strategie neu überschrieben.

\begin{figure}[H]
\centering
\begin{tikzpicture}[
    strat/.style={draw, rounded corners, fill=blue!10, minimum width=4.5cm, minimum height=1cm, font=\small},
    plan/.style={draw, rounded corners, fill=yellow!20, minimum width=3cm, minimum height=0.8cm, font=\small\ttfamily},
    arrow/.style={-Stealth, thick},
    write/.style={-Stealth, thick, red!70!black},
    read/.style={-Stealth, thick, green!50!black, dashed},
]
\node[plan] (plan) at (5,0) {BahnhofPlan};

\node[strat] (s1) at (0, 2.5) {EinfacheFahrtStrategie};
\node[strat] (s2) at (0, 0)   {EinseitigesWartenStrategie};
\node[strat] (s3) at (0,-2.5) {BeidseitigesWartenStrategie};

\draw[write] (s1.east) -- node[above, font=\tiny, red!70!black] {schreibt Hin+Rück} (plan.west |- s1);
\draw[write] (s2.east) -- node[above, font=\tiny, red!70!black] {schreibt Rück+Wa} (plan.west);
\draw[write] (s3.east) -- node[below, font=\tiny, red!70!black] {schreibt Rück+Wa} (plan.west |- s3);

\draw[read, bend left=15] (plan.south west) to node[below right, font=\tiny, green!50!black] {liest ankunfthin} (s2.south east);
\draw[read, bend right=20] (plan.south) to node[below, font=\tiny, green!50!black] {liest ankunfthin} (s3.east);
\end{tikzpicture}
\caption{Alle drei Strategien teilen sich dasselbe \texttt{BahnhofPlan}-Objekt. Strategien 2 und 3 lesen die Hinfahrtdaten, die Strategie 1 hinterlassen hat.}
\end{figure}

% ─────────────────────────────────────────
\section{Entwurfsänderungen gegenüber dem ursprünglichen Design}

Während der Implementierung wurden gegenüber dem initialen Entwurf mehrere
Designentscheidungen revidiert. Dieses Kapitel dokumentiert die wichtigsten
Änderungen, deren Ursachen und die jeweils gewählte Lösung.

% ── 3.1 char → String ──────────────────────────────────────────────────────
\subsection{\texttt{BahnhofNode.name}: \texttt{char} \textrightarrow{} \texttt{String}}

\textbf{Ursprünglicher Entwurf:} Das Feld \texttt{name} in \texttt{BahnhofNode} war als
\texttt{char} geplant – ein einzelnes Zeichen für die Stationsbezeichnung.

\textbf{Problem:} Mit \texttt{char} lassen sich maximal 26 lateinische Großbuchstaben
(A–Z) oder bei doppelter Kodierung maximal 35 Stationen darstellen. Der
Stresstest (\texttt{13\_stress\_test.txt}) mit 52 Stationen erfordert
zweistellige Namen wie \texttt{AA}, \texttt{AB}, \ldots

\textbf{Lösung:} \texttt{char name} wurde durch \texttt{String name} ersetzt.
\texttt{String} erlaubt beliebig lange Bezeichner ohne sonstigen Codeaufwand, da
alle Formatierungen ohnehin über \texttt{String}-Verkettung erfolgen.

\begin{center}
\begin{tabular}{lll}
\toprule
\textbf{Aspekt} & \textbf{Alt (\texttt{char})} & \textbf{Neu (\texttt{String})} \\
\midrule
Typ             & \texttt{private char name}   & \texttt{private String name}   \\
Max. Stationen  & 26 (A–Z)                     & unbegrenzt                     \\
Vergleich       & \texttt{name == 'A'}         & \texttt{name.equals("A")}      \\
\bottomrule
\end{tabular}
\end{center}

% ── 3.2 Kette + BahnhofNode → BahnhofNode ──────────────────────────────────
\subsection{Konsolidierung von \texttt{Kette} und \texttt{BahnhofNode}}

\textbf{Ursprünglicher Entwurf:} Es gab zwei separate Klassen:

\begin{itemize}
  \item \texttt{BahnhofNode} – enthielt ausschließlich die Zeitdaten eines Bahnhofs
        (\texttt{ankunftHin}, \texttt{abfahrtHin}, \texttt{wartezeit}).
  \item \texttt{Kette} – hielt die Listenverknüpfung (\texttt{prev}, \texttt{next},
        \texttt{distance}, \texttt{current}) sowie eine eigene
        \texttt{isKollision()}-Methode.
\end{itemize}

Die originale \texttt{Kette.isKollision()} verwendete eine einfache
Ungleichungsprüfung ohne Uhrzeit-Wrap-around:

\begin{lstlisting}[language=Java, caption={Ursprüngliche \texttt{Kette.isKollision()} -- ohne Wrap-around}]
public boolean isKollision() {
    if ( (current.getAbfahrthin() > next.getAbfahrtrueck() &&
           next.getAnkunfthin() + 1 < next.getAbfahrtrueck()) ||
         (current.getAbfahrthin() >= current.getAnkunftrueck() &&
           next.getAnkunfthin() + 1 < current.getAnkunftrueck()) ) {
        return true;
    }
    return false;
}
\end{lstlisting}

\textbf{Problem:} Die Trennung erzeugte eine zirkuläre Abhängigkeit
(\texttt{Kette} hält \texttt{BahnhofNode}-Referenzen, aber der Knoten kennt seine
Nachbarn nicht). Die \texttt{Kette}-Formel berücksichtigte keine Zeitwerte, die
Mitternacht überschreiten (z.\,B. Abfahrt 58, Ankunft 3).

\textbf{Lösung:} Die Felder \texttt{next}, \texttt{prev} und \texttt{distanceToNext}
wurden direkt in \texttt{BahnhofNode} verschoben; die Klasse \texttt{Kette} wurde
entfernt. Die neue \texttt{isKollision()} in \texttt{BahnhofNode} arbeitet
minutengenau mit Modulo-60-Wrap-around (siehe Kapitel~\ref{sec:kollision}).

\begin{center}
\begin{tabular}{lll}
\toprule
\textbf{Aspekt}           & \textbf{Alt (zwei Klassen)}          & \textbf{Neu (eine Klasse)}            \\
\midrule
Listenstruktur            & \texttt{Kette.prev/next/distance}    & \texttt{BahnhofNode.prev/next/distanceToNext} \\
Kollisionsprüfung         & \texttt{Kette.isKollision()}         & \texttt{BahnhofNode.isKollision()}    \\
Wrap-around               & Nein                                 & Ja (Modulo 60)                        \\
Traversierung             & über \texttt{Kette}-Objekte          & direkt über \texttt{node.next}        \\
\bottomrule
\end{tabular}
\end{center}

% ── 3.3 Simulation-Klasse entfernt ──────────────────────────────────────────
\subsection{Entfernung der geplanten \texttt{Simulation}-Klasse}

\textbf{Ursprünglicher Entwurf:} Eine Klasse \texttt{Simulation} sollte die
Schleife über alle drei Strategien und die Ausgabesteuerung übernehmen.

\textbf{Entscheidung:} Die \texttt{Simulation}-Klasse wurde nicht implementiert.
Stattdessen übernimmt \texttt{Main.main()} diese Orchestrierung direkt:

\begin{lstlisting}[language=Java, caption={Orchestrierung in \texttt{Main.main()} statt \texttt{Simulation}-Klasse}]
FahrplanStrategie[] strategien = {
    new EinfacheFahrtStrategie(),
    new EinseitigesWartenStrategie(),
    new BeidseitigesWartenStrategie()
};
for (FahrplanStrategie strategie : strategien) {
    daten.plan().reset();
    strategie.berechneFahrplan(daten);
    ausgabeManager.druckeErgebnis(strategie, daten);
}
\end{lstlisting}

\textbf{Begründung:} Die Orchestrierungslogik ist für drei feste Strategien trivial
(eine \texttt{for}-Schleife). Eine eigene Klasse hätte nur Boilerplate erzeugt und die
Wartbarkeit nicht verbessert. \texttt{Main} verbleibt schlank – alle
Berechnungslogik liegt in den Strategie-Klassen.

% ── 3.4 EingabeLeser hinzugefügt ────────────────────────────────────────────
\subsection{Hinzufügen von \texttt{EingabeLeser}}

\textbf{Ursprünglicher Entwurf:} Das Einlesen der Eingabedatei war im ursprünglichen
Entwurf keiner eigenen Klasse zugeordnet; es war implizit Teil von \texttt{Main}
oder einer Hilfsklasse.

\textbf{Lösung:} Die Klasse \texttt{EingabeLeser} wurde als separate I/O-Schicht im
Paket \texttt{de.matse.ihk.io} ergänzt. Sie enthält die statische Methode
\texttt{leseDatei(String dateiPfad)}, die eine Eingabedatei zeilenweise parst und
ein \texttt{SimulationsDaten}-Record zurückgibt.

\textbf{Begründung:} Trennung von Eingabe-I/O und Berechnungslogik folgt dem
\textit{Single Responsibility Principle}. \texttt{Main} bleibt von
Dateilese-Details unberührt; Änderungen am Dateiformat erfordern nur
Änderungen in \texttt{EingabeLeser}.

\begin{center}
\begin{tabular}{lll}
\toprule
\textbf{Aspekt}       & \textbf{Alt}              & \textbf{Neu}                     \\
\midrule
I/O-Zuständigkeit     & implizit in \texttt{Main} & \texttt{EingabeLeser.leseDatei()} \\
Rückgabetyp           & --                        & \texttt{SimulationsDaten} (Record) \\
Paket                 & --                        & \texttt{de.matse.ihk.io}          \\
\bottomrule
\end{tabular}
\end{center}

% ── 3.5 Traversierung in FahrplanStrategie ──────────────────────────────────
\subsection{Traversierung und Wartezeitberechnung in \texttt{FahrplanStrategie}}

\textbf{Ursprünglicher Entwurf:} Die Traversierungslogik (Durchlaufen der
Knotenkette für Hin- und Rückfahrt, Berechnung der Wartezeiten) war im Datenmodell
oder einer \texttt{Simulation}-Klasse vorgesehen.

\textbf{Lösung:} Die gesamte Traversierungslogik wurde als \textbf{statische Methoden}
in der abstrakten Klasse \texttt{FahrplanStrategie} implementiert:

\begin{center}
\begin{tabular}{ll}
\toprule
\textbf{Methode}                  & \textbf{Beschreibung}                          \\
\midrule
\texttt{hinfahrtOhneKollision()}  & Vorwärts-Traversierung, keine Wartezeit        \\
\texttt{rueckfahrtOhneKollision()} & Rückwärts-Traversierung, keine Wartezeit      \\
\texttt{hinfahrtMitKollision()}   & Vorwärts mit Wartezeit bei Kollision           \\
\texttt{rueckfahrtMitKollision()} & Rückwärts mit Wartezeit bei Kollision          \\
\texttt{beidseitigesWarten()}     & Gleichmäßige Verteilung der Wartezeit Hin+Rück \\
\texttt{warteZeitHinReset()}      & Setzt alle \texttt{wartezeitHin} auf null zurück \\
\texttt{warteZeitRueckReset()}    & Setzt alle \texttt{wartezeitRueck} auf null zurück \\
\bottomrule
\end{tabular}
\end{center}

\textbf{Begründung:} Statische Hilfsmethoden in der Basisklasse ermöglichen Code-Reuse
zwischen allen drei Strategien, ohne den Datenknoten mit Algorithmuslogik zu
belasten (\textit{Separation of Concerns}). Die konkrete Strategie wählt, welche
Kombination von Hilfsmethoden sie aufruft.

% ── 3.6 wartezeit → wartezeitHin + wartezeitRueck ───────────────────────────
\subsection{Umbenennung: \texttt{wartezeit} \textrightarrow{} \texttt{wartezeitHin} + \texttt{wartezeitRueck}}

\textbf{Ursprünglicher Entwurf:} \texttt{BahnhofNode} hatte ein einziges Feld
\texttt{wartezeit}, das undifferenziert die Wartezeit speicherte.

\textbf{Problem:} Bei beidseitigem Warten müssen Hin- und Rückfahrt-Wartezeiten
unabhängig voneinander verwaltet werden. Ein einzelnes Feld würde Mehrdeutigkeit
erzeugen.

\textbf{Lösung:} Das Feld wurde in zwei separate Felder aufgeteilt:
\begin{itemize}
  \item \texttt{Integer wartezeitHin} – Wartezeit der Hinfahrt an diesem Knoten
  \item \texttt{Integer wartezeitRueck} – Wartezeit der Rückfahrt an diesem Knoten
\end{itemize}

Der ursprüngliche Name ist noch im auskommentierten \texttt{toString()} von
\texttt{BahnhofNode} als historischer Überrest sichtbar.

\begin{center}
\begin{tabular}{lll}
\toprule
\textbf{Aspekt}       & \textbf{Alt}                   & \textbf{Neu}                                      \\
\midrule
Feld(er)              & \texttt{Integer wartezeit}     & \texttt{Integer wartezeitHin, wartezeitRueck}     \\
Strafe EinseitigesWarten & $(\sum w)^2$                & $(\sum \mathtt{wartezeitRueck})^2$                \\
Strafe BeidseitigesWarten & --                         & $(\sum \mathtt{wartezeitHin})^2 + (\sum \mathtt{wartezeitRueck})^2$ \\
\bottomrule
\end{tabular}
\end{center}

% ── 3.7 Verbesserte isKollision ──────────────────────────────────────────────
\subsection{Verbesserung der Kollisionserkennung}

Die ursprüngliche \texttt{Kette.isKollision()} (Abschnitt~3.2) verwendete einfache
Ungleichungsvergleiche und erkannte keine Kollisionen, wenn Zeitintervalle
Mitternacht (Minute 59 $\to$ 0) überspannen.

Die neue \texttt{BahnhofNode.isKollision()} iteriert minutengenau alle 60 möglichen
Zeitpunkte und prüft über \texttt{isTimeInInterval()} mit Modulo-60-Logik, ob ein
Zug zu diesem Zeitpunkt auf dem Segment ist. Kollision liegt vor, wenn beide Züge
zur gleichen Minute auf demselben Streckenabschnitt fahren.

\begin{center}
\begin{tabular}{lll}
\toprule
\textbf{Aspekt}            & \textbf{Alt (\texttt{Kette})}      & \textbf{Neu (\texttt{BahnhofNode})}        \\
\midrule
Methode                    & \texttt{Kette.isKollision()}       & \texttt{BahnhofNode.isKollision()}         \\
Algorithmus                & Zwei Ungleichungen                 & Minutenweise Überlappung (60 Iterationen)  \\
Wrap-around                & Nicht berücksichtigt               & \texttt{isTimeInInterval()} mit \% 60      \\
Robustheit Mitternacht     & Fehlerhafte Ergebnisse möglich     & Korrekt für alle Zeitwerte 0–59            \\
\bottomrule
\end{tabular}
\end{center}

% ── Übersichtstabelle ────────────────────────────────────────────────────────
\subsection*{Übersicht aller Entwurfsänderungen}

\begin{center}
\begin{tabular}{clll}
\toprule
\textbf{Nr.} & \textbf{Änderung}                     & \textbf{Ursache}                        & \textbf{Auswirkung}             \\
\midrule
1 & \texttt{char} $\to$ \texttt{String}          & Stationsanzahl $>$ 26                   & Beliebig lange Namen möglich    \\
2 & \texttt{Kette} entfernt                       & Zirkuläre Abhängigkeit                  & \texttt{next/prev} in Node      \\
3 & \texttt{Simulation} nicht implementiert       & Triviale Logik                          & Orchestrierung in \texttt{Main} \\
4 & \texttt{EingabeLeser} hinzugefügt             & SRP / I/O-Trennung                      & Klare Schichtentrennung         \\
5 & Traversierung in \texttt{FahrplanStrategie}   & Separation of Concerns                  & Strategien nutzen gemeinsame Helpers \\
6 & \texttt{wartezeit} $\to$ 2 Felder             & Beidseitiges Warten erfordert Trennung  & Genaue Strafenberechnung        \\
7 & \texttt{isKollision()} verbessert             & Wrap-around-Fehler                      & Korrekte Mitternachtserkennung  \\
\bottomrule
\end{tabular}
\end{center}

% ─────────────────────────────────────────
\section{Einlesen und Initialisieren der Daten}

\subsection{Eingabedateiformat}

Die Eingabedateien haben folgendes Format (Schlüsselwort in einer Zeile, Wert in der nächsten):

\begin{verbatim}
Strecke:
A B C D E F G
Abstaende:
5 8 6 7 3 2
Start Hinfahrt:
17
\end{verbatim}

\begin{itemize}
    \item \textbf{Strecke}: Die Bahnhofsnamen, durch Leerzeichen getrennt.
    \item \textbf{Abstaende}: Die Fahrtzeit (in Minuten) zwischen je zwei aufeinanderfolgenden Bahnhöfen. Bei $n$ Bahnhöfen gibt es $n-1$ Abstände.
    \item \textbf{Start Hinfahrt}: Die Abfahrtszeit des Zuges am ersten Bahnhof (Minuten, Modulo 60).
\end{itemize}

\subsection{Einleseprozess (EingabeLeser)}

Die Klasse \texttt{EingabeLeser} liest die Datei zeilenweise ein und erkennt die Schlüsselwörter per einfachem String-Vergleich:

\begin{lstlisting}[caption={Kernteil von EingabeLeser.leseDatei()}]
List<String> zeilen = Files.readAllLines(Path.of(dateiPfad));
BahnhofPlan bahnhofPlan = new BahnhofPlan();
int startZeit = 0;

for (int i = 0; i < zeilen.size(); i++) {
    String aktuelleZeile = zeilen.get(i).trim();
    if (aktuelleZeile.equals("Strecke:")) {
        String[] bahnhoefe = zeilen.get(i + 1).split("\\s+");
        for (String name : bahnhoefe)
            bahnhofPlan.addBahnhof(name);
    } else if (aktuelleZeile.equals("Abstaende:")) {
        String[] abstaendeStr = zeilen.get(i + 1).trim().split("\\s+");
        int[] abstaende = new int[abstaendeStr.length];
        for (int j = 0; j < abstaendeStr.length; j++)
            abstaende[j] = Integer.parseInt(abstaendeStr[j]);
        bahnhofPlan.setAbstaende(abstaende);
    } else if (aktuelleZeile.equals("Start Hinfahrt:")) {
        startZeit = Integer.parseInt(zeilen.get(i + 1).trim());
    }
}
return new SimulationsDaten(bahnhofPlan, startZeit);
\end{lstlisting}

\begin{enumerate}
    \item Beim Schlüsselwort \texttt{Strecke:} wird für jeden Bahnhofsnamen ein neuer \texttt{BahnhofNode} angelegt und per \texttt{addBahnhof()} an die Liste angehängt.
    \item Beim Schlüsselwort \texttt{Abstaende:} werden die Abstände als \texttt{int[]} geparst und per \texttt{setAbstaende()} jedem Knoten sein \texttt{distanceToNext}-Wert zugewiesen.
    \item Beim Schlüsselwort \texttt{Start Hinfahrt:} wird die Startzeit als \texttt{int} gespeichert.
    \item Das Ergebnis wird als \texttt{SimulationsDaten}-Record zurückgegeben.
\end{enumerate}

% ─────────────────────────────────────────
\section{Kollisionserkennung}

\subsection{Definition einer Kollision}

Eine \textbf{Kollision} tritt auf, wenn sich Hin- und Rückzug \textbf{gleichzeitig auf demselben Streckenabschnitt} befinden. Ein Streckenabschnitt liegt zwischen zwei aufeinanderfolgenden Bahnhöfen, z.\,B. zwischen Bahnhof B und Bahnhof C.

\begin{itemize}
    \item \textbf{Hinfahrt} auf dem Abschnitt B$\to$C: Der Zeitraum beginnt mit der Abfahrt des Hinzuges bei B (\texttt{abfahrthin}) und endet mit seiner Ankunft bei C (\texttt{ankunfthin}).
    \item \textbf{Rückfahrt} auf dem Abschnitt C$\to$B: Der Zeitraum beginnt mit der Abfahrt des Rückzuges bei C (\texttt{abfahrtrueck}) und endet mit seiner Ankunft bei B (\texttt{ankunftrueck}).
\end{itemize}

Eine Kollision liegt vor, genau dann wenn sich diese zwei Zeitintervalle \textbf{überschneiden}.

\subsection{Uhrzeitmodell (Modulo-60-Uhr)}

Die Zeit wird als Minutenwert modulo 60 dargestellt, d.\,h. nach Minute 59 folgt wieder Minute 0. Ein Intervall kann daher über den „Stundenwechsel" gehen (z.\,B. von Minute 55 bis Minute 05).

\subsection{Implementierung: isKollision()}

Die Methode \texttt{isKollision()} gehört zur Klasse \texttt{BahnhofNode} und prüft den Abschnitt \emph{zwischen diesem Knoten und seinem Nachfolger}:

\begin{lstlisting}[caption={BahnhofNode.isKollision()}]
public boolean isKollision() {
    if (this.next == null) return false; // letzter Bahnhof: kein Abschnitt danach

    int hStart = this.abfahrthin;            // Hinzug faehrt ab
    int hEnde  = this.next.getAnkunfthin();  // Hinzug kommt an

    int rStart = this.next.getAbfahrtrueck(); // Rueckzug faehrt ab
    int rEnde  = this.ankunftrueck;           // Rueckzug kommt an

    return checkOverlap(hStart, hEnde, rStart, rEnde);
}
\end{lstlisting}

\begin{lstlisting}[caption={Intervall-Überschneidungsprüfung}]
private boolean checkOverlap(int s1, int e1, int s2, int e2) {
    for (int m = 0; m < 60; m++) {
        if (isTimeInInterval(m, s1, e1) && isTimeInInterval(m, s2, e2))
            return true;
    }
    return false;
}

private boolean isTimeInInterval(int t, int start, int ende) {
    if (start <= ende) {
        return t >= start && t <= ende;   // Normalfall
    } else {
        return t >= start || t <= ende;   // Uebermitternacht
    }
}
\end{lstlisting}

Die Methode \texttt{checkOverlap()} iteriert alle 60 möglichen Minutenwerte und prüft, ob ein Zeitpunkt \texttt{m} in \textbf{beiden} Intervallen liegt. Die Hilfsmethode \texttt{isTimeInInterval()} behandelt dabei korrekt den Sonderfall, dass ein Intervall über den Stundenwechsel geht.

\subsection{Konkretes Beispiel: Kollisionsprüfung an Bahnhof B (Beispiel2)}

\begin{beispielbox}[title={Beispiel2: Strecke A--B--C--D--E--F--G, Abstände: 5 8 6 7 3 2, Start: 17}]
Hinfahrt (ohne Kollisionsbehandlung):
\begin{center}
\begin{tabular}{lccccccc}
\toprule
 & A & B & C & D & E & F & G \\
\midrule
Ankunft Hin & -- & 22 & 31 & 38 & 46 & 50 & 53 \\
Abfahrt Hin & 17 & 23 & 32 & 39 & 47 & 51 & -- \\
\bottomrule
\end{tabular}
\end{center}

Rückfahrt (ohne Kollisionsbehandlung, Startzeit = 54):
\begin{center}
\begin{tabular}{lccccccc}
\toprule
 & A & B & C & D & E & F & G \\
\midrule
Abfahrt Rück & -- & 25 & 16 & 09 & 01 & 57 & 54 \\
Ankunft Rück & 30 & 24 & 15 & 08 & 00 & 56 & -- \\
\bottomrule
\end{tabular}
\end{center}

\textbf{Kollisionsprüfung am Abschnitt B$\to$C} (Knoten B, \texttt{isKollision()}):
\begin{align*}
\text{Hinzug auf B}\to\text{C:} &\quad [\texttt{abfahrthin}(B),\; \texttt{ankunfthin}(C)] = [23,\; 31] \\
\text{Rückzug auf C}\to\text{B:} &\quad [\texttt{abfahrtrueck}(C),\; \texttt{ankunftrueck}(B)] = [16,\; 24]
\end{align*}

Schnittmenge: $t = 23$ und $t = 24$ liegen in \textbf{beiden} Intervallen $\Rightarrow$ \textbf{KOLLISION!}

Die Ausgabe markiert diesen Abschnitt mit einem \texttt{x} zwischen B und C.
\end{beispielbox}

% ─────────────────────────────────────────
\section{Strategie 1: Einfache Fahrt}

\subsection{Idee und Zweck}

Die \textbf{Einfache Fahrt} ist die einfachste aller drei Strategien: Sie berechnet die schnellstmögliche Hin- und Rückfahrt, als ob die Strecke \emph{zweigleisig} wäre. Kollisionen werden dabei vollständig ignoriert. Die Strategie dient als \textbf{Referenzgrundlage} -- sie liefert die minimale Fahrtdauer ohne jede Wartezeit und zeigt durch die Kollisionsmarkierung (\texttt{x}), wo bei echtem Eingleis-Betrieb Probleme entstünden.

Konkret bedeutet das: Die Hinfahrt fährt durch, ohne auf den Rückzug zu warten. Der Rückzug startet unmittelbar nach Ankunft des Hinzuges am Zielbahnhof und fährt ebenfalls durch, ohne auf den Hinzug zu warten. Begegnungen auf freier Strecke werden in dieser Strategie nicht als Problem behandelt.

\subsection{Verbale Beschreibung des Algorithmus}

\textbf{Warum +1 Minute Haltezeit?}

An jedem Zwischenbahnhof hält der Zug genau eine Minute (Tür auf, Tür zu). Deshalb gilt: Die Abfahrtszeit an einem Bahnhof ist immer genau eine Minute nach der Ankunftszeit. Nur am letzten Bahnhof gibt es keine Abfahrt mehr.

\textbf{Warum Modulo 60?}

Die Uhr läuft im 60-Minuten-Takt. Nach Minute 59 folgt Minute 0. Die Modulo-60-Rechnung stellt sicher, dass Zeiten immer im Bereich 0--59 bleiben, auch wenn ein Zug über den Stundenwechsel hinausfährt.

\textbf{Hinfahrt -- Schritt für Schritt}:

\begin{enumerate}
    \item Dem ersten Bahnhof (Head) wird die Startzeit als Abfahrtszeit zugewiesen: \texttt{abfahrthin = startZeit}.
    \item Für jeden folgenden Bahnhof wird die Ankunftszeit berechnet: Die Fahrtzeit (\texttt{distanceToNext}) des vorherigen Bahnhofs wird zur Abfahrtszeit addiert -- modulo 60. Das Ergebnis ist die \textbf{Ankunftszeit} am nächsten Bahnhof.
    \item Die \textbf{Abfahrtszeit} des nächsten Bahnhofs ist die Ankunftszeit plus eine Minute Haltezeit.
    \item Dies wird für jeden Bahnhof wiederholt, bis der letzte Bahnhof (Tail) erreicht ist. Der Tail hat keine Abfahrtszeit (er ist das Ziel).
\end{enumerate}

\textbf{Rückfahrt -- Schritt für Schritt}:

\begin{enumerate}
    \item Der Rückzug startet eine Minute nach Ankunft des Hinzuges am Zielbahnhof: \texttt{startZeitRueck = (ankunfthin(tail) + 1) \% 60}.
    \item Die Berechnung läuft rückwärts: vom letzten Bahnhof (Tail) zum ersten (Head). Dafür wird die \texttt{prev}-Zeiger-Kette der doppelt verketteten Liste genutzt.
    \item Für jeden Vorgängerknoten gilt dieselbe Logik wie bei der Hinfahrt: Ankunft = vorherige Abfahrt + Fahrzeit, Abfahrt = Ankunft + 1 Minute.
    \item Der erste Bahnhof (Head) erhält am Ende nur eine Ankunftszeit, aber keine Abfahrtszeit (er ist das Ziel der Rückfahrt).
\end{enumerate}

\textbf{Kollisionserkennung in der Ausgabe}:

Nach Berechnung beider Fahrten wird für jeden Abschnitt \texttt{isKollision()} aufgerufen. Wenn ein Abschnitt kollisionsbehaftet ist, erscheint ein \texttt{x} in der Ausgabe zwischen den betroffenen Bahnhöfen. Die Kollision wird jedoch \textbf{nicht behoben} -- die Zeitwerte bleiben unverändert.

\subsection{Formeln}

\begin{align*}
    \texttt{ankunfthin}_{i+1} &= (\texttt{abfahrthin}_i + \texttt{distanceToNext}_i) \bmod 60 \\
    \texttt{abfahrthin}_{i+1} &= (\texttt{ankunfthin}_{i+1} + 1) \bmod 60 \\[6pt]
    \texttt{startZeitRueck}   &= (\texttt{ankunfthin}_\text{tail} + 1) \bmod 60 \\
    \texttt{ankunftrueck}_{i-1} &= (\texttt{abfahrtrueck}_i + \texttt{distanceToNext}_{i-1}) \bmod 60 \\
    \texttt{abfahrtrueck}_{i-1} &= (\texttt{ankunftrueck}_{i-1} + 1) \bmod 60
\end{align*}

\textbf{Strafe}: Da keine Wartezeiten anfallen, ist die Strafe immer $0$ (auch wenn Kollisionen vorhanden sind).

\subsection{Schrittweise Berechnung am Beispiel2}

\begin{beispielbox}[title={Einfache Fahrt -- Beispiel2: A--B--C--D--E--F--G, Abstände: 5 8 6 7 3 2, Start: 17}]

\textbf{Hinfahrt} -- Berechnung Knoten für Knoten:

\begin{center}
\begin{tabular}{c|c|c|p{7cm}}
\toprule
\textbf{Knoten} & \textbf{An} & \textbf{Ab} & \textbf{Berechnung} \\
\midrule
A & -- & 17 & Startzeit direkt gesetzt \\
B & $(17+5)\bmod60 = 22$ & $22+1 = 23$ & Abstand A$\to$B = 5 \\
C & $(23+8)\bmod60 = 31$ & $31+1 = 32$ & Abstand B$\to$C = 8 \\
D & $(32+6)\bmod60 = 38$ & $38+1 = 39$ & Abstand C$\to$D = 6 \\
E & $(39+7)\bmod60 = 46$ & $46+1 = 47$ & Abstand D$\to$E = 7 \\
F & $(47+3)\bmod60 = 50$ & $50+1 = 51$ & Abstand E$\to$F = 3 \\
G & $(51+2)\bmod60 = 53$ & -- & Ziel, keine Abfahrt \\
\bottomrule
\end{tabular}
\end{center}

\textbf{Rückfahrt} -- Startzeit: $(53+1)\bmod60 = 54$, Berechnung rückwärts:

\begin{center}
\begin{tabular}{c|c|c|p{7cm}}
\toprule
\textbf{Knoten} & \textbf{Ab} & \textbf{An} & \textbf{Berechnung} \\
\midrule
G & 54 & -- & Startzeit direkt gesetzt \\
F & $56+1 = 57$ & $(54+2)\bmod60 = 56$ & Abstand F$\to$G = 2 \\
E & $0+1 = 1$ & $(57+3)\bmod60 = 0$ & Abstand E$\to$F = 3 \\
D & $8+1 = 9$ & $(1+7)\bmod60 = 8$ & Abstand D$\to$E = 7 \\
C & $15+1 = 16$ & $(9+6)\bmod60 = 15$ & Abstand C$\to$D = 6 \\
B & $24+1 = 25$ & $(16+8)\bmod60 = 24$ & Abstand B$\to$C = 8 \\
A & -- & $(25+5)\bmod60 = 30$ & Ziel, keine Abfahrt \\
\bottomrule
\end{tabular}
\end{center}

\textbf{Ergebnis}:
\begin{center}
\begin{tabular}{r|ccccccc}
 & A & B & C & D & E & F & G \\
\hline
An Hin  &    & 22 & 31 & 38 & 46 & 50 & 53 \\
Wa Hin  &    &    &    &    &    &    &    \\
Ab Hin  & 17 & 23 & 32 & 39 & 47 & 51 &    \\
\hline
Bahnhof & A  & B\textbf{x} & C & D & E & F & G \\
\hline
Ab Rück &    & 25 & 16 & 09 & 01 & 57 & 54 \\
Wa Rück &    &    &    &    &    &    &    \\
An Rück & 30 & 24 & 15 & 08 & 00 & 56 &    \\
\end{tabular}
\end{center}

Gesamtdauer Hin/Rück: 36 / 36 \quad Wartezeit: 0 / 0 \quad Strafe: 0 \\[4pt]
Das \textbf{x} bei B zeigt: Auf dem Abschnitt B$\to$C kollidierten die Züge. Dies bleibt unbehandelt.
\end{beispielbox}

\begin{hinweisbox}
Die Hinfahrtdaten (\texttt{ankunfthin}, \texttt{abfahrthin}), die \texttt{EinfacheFahrtStrategie} in die Knoten schreibt, werden von den Strategien~2 und~3 \emph{unverändert weitergenutzt} (siehe Abschnitt~\ref{sec:referenz}). Die Strategien berechnen die Hinfahrt \emph{nicht} neu.
\end{hinweisbox}

% ─────────────────────────────────────────
\section{Strategie 2: Einseitiges Warten}

\subsection{Idee und Zweck}

Beim \textbf{Einseitigen Warten} löst das Programm Kollisionen auf, indem es den \textbf{Rückzug an bestimmten Bahnhöfen warten lässt}. Der Hinzug fährt unverändert durch (identisch zur Einfachen Fahrt). Der Rückzug muss an einer Kollisionsstelle so lange warten, bis der Hinzug den gemeinsamen Streckenabschnitt vollständig verlassen hat -- erst dann darf der Rückzug in diesen Abschnitt einfahren.

Dieses Verfahren garantiert, dass \textbf{keine Kollision mehr entsteht}. Der Nachteil: Der Rückzug trägt die gesamte Verzögerung allein. Je länger er wartet, desto größer ist die quadratische Strafe.

\subsection{Verbale Beschreibung des Algorithmus}

Die Methode \texttt{rueckfahrtMitKollision()} arbeitet in \textbf{zwei Phasen}:

\textbf{Phase 1 -- Ausgangsbasis ohne Kollision}:

Zunächst wird die Rückfahrt exakt wie bei der Einfachen Fahrt berechnet (Methode \texttt{rueckfahrtOhneKollision()}). Das Ergebnis ist noch nicht kollisionsfrei, dient aber als Startpunkt für die Kollisionskorrektur in Phase 2.

\textbf{Phase 2 -- Rückwärtige Traversierung mit Kollisionskorrektur}:

Der Algorithmus läuft vom vorletzten Bahnhof (tail.prev) rückwärts bis zum zweiten Bahnhof (head.next). Für jeden Knoten $i$ in dieser Reihenfolge werden folgende Schritte ausgeführt:

\begin{enumerate}
    \item \textbf{Zeiten neu berechnen}: Da ein vorheriger Schritt möglicherweise Wartezeiten eingebaut hat, die alle nachfolgenden Abfahrtszeiten verschieben, werden die Ankunfts- und Abfahrtszeiten des aktuellen Knotens aus der bereits berechneten Abfahrtszeit des \emph{nächsten} Knotens neu abgeleitet.
    \item \textbf{Kollision am Vorgänger prüfen}: Es wird geprüft, ob der \emph{Vorgänger-Knoten} ($i-1$) auf seinem Abschnitt nach $i$ eine Kollision hat. Der Vorgänger wird geprüft, weil der Rückzug dort in den gefährdeten Abschnitt \emph{einfährt} -- dort muss er warten.
    \item \textbf{Wartezeit berechnen}: Wenn eine Kollision erkannt wurde, wird die nötige Wartezeit des Rückzuges bei Knoten $i$ berechnet. Der Rückzug muss mindestens so lange warten, bis der Hinzug bei Knoten $i$ ankommt -- denn erst dann hat der Hinzug den Abschnitt zwischen $i-1$ und $i$ verlassen.

    Die Wartezeit ergibt sich als Differenz zwischen der Ankunftszeit des Hinzuges und der Ankunftszeit des Rückzuges an Knoten $i$ (plus Modulo-60-Korrektur für den Stundenwechsel):
    \begin{align*}
        \texttt{wartezeitRueck}_i = (\texttt{ankunfthin}_i - \texttt{ankunftrueck}_i + 60) \bmod 60
    \end{align*}

    \item \textbf{Abfahrtszeit anpassen}: Die Abfahrtszeit des Rückzuges bei Knoten $i$ wird nach hinten verschoben:
    \begin{align*}
        \texttt{abfahrtrueck}_i = (\texttt{ankunftrueck}_i + \texttt{wartezeitRueck}_i + 1) \bmod 60
    \end{align*}
    Der $+1$-Term ist wiederum die Haltezeit.
\end{enumerate}

\textbf{Warum rückwärts und nicht vorwärts?}

Der Algorithmus läuft vom letzten Bahnhof in Richtung erster Bahnhof, weil der Rückzug ebenfalls in dieser Richtung fährt. Wenn bei Bahnhof $i$ eine Wartezeit eingefügt wird, verschiebt sich die Abfahrtszeit dort nach hinten -- das wirkt sich auf alle \emph{vorher} besuchten Bahnhöfe (also weiter vorne in der Strecke) aus. Indem wir rückwärts (in Fahrtrichtung des Rückzuges vorwärts) laufen, werden diese Kaskaden korrekt aufgefangen.

\textbf{Warum wird der Vorgänger-Knoten auf Kollision geprüft, nicht der aktuelle?}

\texttt{isKollision()} ist am Knoten definiert und prüft den Abschnitt \emph{zwischen diesem Knoten und seinem Nachfolger} (in Richtung der Hinfahrt). Der Rückzug fährt auf diesem Abschnitt von $i$ nach $i-1$ -- also von rechts nach links in der Liste. Der für den Rückzug gefährliche Abschnitt zwischen $i-1$ und $i$ wird durch \texttt{isKollision()} am Knoten $i-1$ geprüft: \texttt{current.getPrev().isKollision()}.

\textbf{Strafe}:
\begin{align*}
    \text{Strafe} = \left(\sum_i \texttt{wartezeitRueck}_i\right)^2
\end{align*}

\subsection{Lückenlose Schritt-für-Schritt-Berechnung: Beispiel2}

\begin{beispielbox}[title={Einseitiges Warten -- Beispiel2: A--B--C--D--E--F--G, Abstände: 5 8 6 7 3 2, Start: 17}]

Die Hinfahrt ist identisch zur Einfachen Fahrt und bereits in die Knoten geschrieben (von Strategie~1):

\begin{center}
\begin{tabular}{r|ccccccc}
 & A & B & C & D & E & F & G \\
\hline
An Hin & -- & 22 & 31 & 38 & 46 & 50 & 53 \\
Ab Hin & 17 & 23 & 32 & 39 & 47 & 51 & -- \\
\end{tabular}
\end{center}

\textbf{Phase 1} -- Rückfahrt ohne Kollision (Startzeit G = $(53+1)\bmod60 = 54$):

\begin{center}
\begin{tabular}{c|cc|p{7cm}}
\toprule
\textbf{Knoten} & \textbf{Ab Rück} & \textbf{An Rück} & \textbf{Berechnung} \\
\midrule
G & 54 & -- & Startzeit \\
F & 57 & $(54+2)=56$ & Abstand F$\to$G = 2 \\
E & 1 & $(57+3)\bmod60=0$ & Abstand E$\to$F = 3 \\
D & 9 & $(1+7)=8$ & Abstand D$\to$E = 7 \\
C & 16 & $(9+6)=15$ & Abstand C$\to$D = 6 \\
B & 25 & $(16+8)=24$ & Abstand B$\to$C = 8 \\
A & -- & $(25+5)=30$ & Ziel der Rückfahrt \\
\bottomrule
\end{tabular}
\end{center}

\textbf{Phase 2} -- Rückwärtstraversierung, \texttt{current} startet bei F:

\medskip
\textit{Iteration 1: current = F}

\begin{itemize}
    \item Zeiten neu: F.An=$(G.Ab+2)\bmod60=(54+2)=56$, F.Ab=$57$
    \item Vorausschau: E.An=$(57+3)=0$
    \item Kollision bei E? (Abschnitt E$\to$F): Hin$[47,50]$, Rück$[57,0]$ $\to$ kein Überlapp $\Rightarrow$ keine Kollision
    \item F.wartezeitRueck = 0, F.Ab bleibt 57
\end{itemize}

\textit{Iteration 2: current = E}

\begin{itemize}
    \item Zeiten neu: E.An=$(57+3)=0$, E.Ab=$1$
    \item Vorausschau: D.An=$(1+7)=8$
    \item Kollision bei D? (Abschnitt D$\to$E): Hin$[39,46]$, Rück$[1,8]$ $\to$ kein Überlapp $\Rightarrow$ keine Kollision
    \item E.wartezeitRueck = 0
\end{itemize}

\textit{Iteration 3: current = D}

\begin{itemize}
    \item Zeiten neu: D.An=$(1+7)=8$, D.Ab=$9$
    \item Vorausschau: C.An=$(9+6)=15$
    \item Kollision bei C? (Abschnitt C$\to$D): Hin$[32,38]$, Rück$[9,15]$ $\to$ kein Überlapp $\Rightarrow$ keine Kollision
    \item D.wartezeitRueck = 0
\end{itemize}

\textit{Iteration 4: current = C}

\begin{itemize}
    \item Zeiten neu: C.An=$(9+6)=15$, C.Ab=$16$
    \item Vorausschau: B.An=$(16+8)=24$
    \item Kollision bei B? (Abschnitt B$\to$C): Hin$[23,31]$, Rück$[16,24]$

    $\to$ $t=23$: in $[23,31]$? Ja. In $[16,24]$? Ja. \textbf{KOLLISION!}
    \item Wartezeit: $\texttt{wartezeitRueck}(C) = (31 - 15 + 60)\bmod60 = 76\bmod60 = 16$ Minuten

    \textit{Erklärung}: Der Hinzug kommt bei C um Min.~31 an. Der Rückzug kommt schon um Min.~15 an. Er muss $31-15=16$ Minuten warten, bis der Hinzug an C vorbei ist.
    \item C.Ab $= (15 + 16 + 1)\bmod60 = 32$
\end{itemize}

\textit{Iteration 5: current = B}

\begin{itemize}
    \item Zeiten neu: B.An=$(C.Ab+8)=(32+8)=40$, B.Ab=$41$
    \item Vorausschau: A.An=$(41+5)=46$
    \item Kollision bei A? (Abschnitt A$\to$B): Hin$[17,22]$, Rück$[41,46]$ $\to$ kein Überlapp $\Rightarrow$ keine Kollision
    \item B.wartezeitRueck = 0
\end{itemize}

\textit{Nach der Schleife}: A.An $= (41+5)\bmod60 = 46$

\textbf{Endergebnis}:
\begin{center}
\begin{tabular}{r|ccccccc}
 & A & B & C & D & E & F & G \\
\hline
An Hin  &    & 22 & 31 & 38 & 46 & 50 & 53 \\
Wa Hin  &    &    &    &    &    &    &    \\
Ab Hin  & 17 & 23 & 32 & 39 & 47 & 51 &    \\
\hline
Bahnhof & A  & B  & C  & D  & E  & F  & G  \\
\hline
Ab Rück &    & 41 & 32 & 09 & 01 & 57 & 54 \\
Wa Rück &    &    & \textbf{16} &    &    &    &    \\
An Rück & 46 & 40 & 15 & 08 & 00 & 56 &    \\
\end{tabular}
\end{center}

Gesamtdauer Hin/Rück: 36 / 52 \quad Wartezeit: 0 / 16 \quad Strafe: $16^2 = \mathbf{256}$

Der Rückzug wartet 16 Minuten bei C, bis der Hinzug C passiert hat (Min.~31). Danach kann er ungehindert von C nach B nach A fahren.
\end{beispielbox}

% ─────────────────────────────────────────
\section{Strategie 3: Beidseitiges Warten}

\subsection{Idee und Zweck}

Das \textbf{Beidseitige Warten} ist die intelligenteste der drei Strategien. Sie behebt Kollisionen, indem sie die entstehende Wartezeit \textbf{möglichst gleichmäßig auf Hin- und Rückzug aufteilt}.

Der Grundgedanke: Wenn zwei Züge auf einer Strecke aufeinandertreffen, müssen zusammen $W$ Minuten Wartezeit erzeugt werden. Statt den gesamten Betrag einem Zug aufzubürden (wie beim Einseitigen Warten), lässt man jeden Zug die Hälfte warten. Das klingt fair -- aber es ist auch \textbf{mathematisch optimal}, weil die Strafe quadratisch berechnet wird:

\begin{align*}
\text{Einseitiges Warten:} & \quad (W)^2 \\
\text{Beidseitiges Warten:} & \quad \left(\frac{W}{2}\right)^2 + \left(\frac{W}{2}\right)^2 = \frac{W^2}{4} + \frac{W^2}{4} = \frac{W^2}{2}
\end{align*}

Das Beidseitige Warten erzielt also \textbf{immer mindestens die halbe Strafe} des Einseitigen Wartens.

Noch besser: Häufig findet die Strategie eine Startzeit für den Rückzug, bei der \emph{gar keine} Wartezeit nötig ist -- dann ist die Strafe 0.

\subsection{Verbale Beschreibung des Algorithmus}

Die Strategie arbeitet in \textbf{zwei Schritten}:

\subsubsection{Schritt 1: Optimale Rückfahrt-Startzeit suchen}

Der Algorithmus probiert \textbf{alle 60 möglichen Startzeiten} (Minute 0 bis 59) für die Rückfahrt durch. Für jede Startzeit wird mit \texttt{rueckfahrtMitKollision()} berechnet, wie die Rückfahrt aussähe (inklusive einseitigem Warten, wo nötig). Die Startzeit, die die \textbf{geringste Strafe} liefert, wird als beste Startzeit gespeichert.

Eine Startzeit wird übersprungen (\texttt{continue}), wenn der vorletzte Bahnhof (direkt neben dem Endbahnhof) noch eine Kollision hat. Das ist ein Sonderfall, bei dem der Algorithmus nicht sauber abschließen kann.

\textbf{Warum alle 60 ausprobieren?} Weil es keine einfache Formel gibt, die direkt die optimale Startzeit berechnet. Die Streckenabstände und Kollisionssituationen sind zu vielfältig. Der Brute-Force-Ansatz ist jedoch sehr effizient: 60 Iterationen $\times$ Anzahl der Knoten -- das läuft in Millisekunden.

\subsubsection{Schritt 2: Wartezeit hälftig aufteilen (falls noch nötig)}

Wenn die beste gefundene Startzeit noch Restkollisionen enthält (d.\,h. der Rückzug musste trotz optimaler Startzeit noch warten), wird die verbleibende Wartezeit hälftig aufgeteilt:

\begin{itemize}
    \item Für jeden Bahnhof mit \texttt{wartezeitRueck > 0} wird die Hälfte dieser Wartezeit dem \textbf{Hinzug} zugewiesen (\texttt{wartezeitHin}).
    \item Die andere Hälfte bleibt beim Rückzug (\texttt{wartezeitRueck}).
    \item Bei ungerader Wartezeit (z.\,B. 7 Min.) bekommt die Hinfahrt 3 Min. und die Rückfahrt 4 Min. -- die Rückfahrt trägt eine Minute mehr.
\end{itemize}

\textbf{Was bedeutet Wartezeit für den Hinzug?}

Wenn der Hinzug an Bahnhof $i$ wartet, fährt er erst später ab. Das verschiebt alle nachfolgenden Ankünfte des Hinzuges nach hinten -- und damit auch den Zeitpunkt, ab dem der Rückzug ohne Kollision einfahren kann. Die Hinfahrt muss deshalb nach der Aufteilung mit \texttt{hinfahrtMitKollision()} neu berechnet werden.

Anschließend wird auch die Rückfahrt mit der angepassten (um die hinzugekommene Hinverzögerung verschobenen) Startzeit neu berechnet.

\subsection{Formeln}

\textbf{Aufteilung der Wartezeit}:
\begin{align*}
\texttt{wartezeitHin}_i &= \left\lfloor \texttt{wartezeitRueck}_i \,/\, 2 \right\rfloor \\
\texttt{wartezeitRueck}_i &= \texttt{wartezeitRueck}_i - \texttt{wartezeitHin}_i
\end{align*}

\textbf{Strafe}:
\begin{align*}
    \text{Strafe} = \left(\sum_i \texttt{wartezeitHin}_i\right)^2 + \left(\sum_i \texttt{wartezeitRueck}_i\right)^2
\end{align*}

\subsection{Lückenlose Berechnung: Beispiel2 (einfacher Fall -- keine Wartezeit nötig)}

\begin{beispielbox}[title={Beidseitiges Warten -- Beispiel2: A--B--C--D--E--F--G, Abstände: 5 8 6 7 3 2, Start: 17}]

\textbf{Schritt 1: Suche der besten Startzeit}

Der Algorithmus probiert alle 60 Startzeiten. Relevant ist die Kollision auf dem Abschnitt B$\to$C. Die Hinfahrt belegt diesen Abschnitt im Intervall $[23, 31]$.

Damit der Rückzug diesen Abschnitt \emph{nicht gleichzeitig} nutzt, muss das Rückfahrtintervall auf C$\to$B entweder \textbf{vor Minute 23} oder \textbf{nach Minute 31} liegen. Bei Startzeit G = 10:

\begin{center}
\begin{tabular}{c|cc|p{6cm}}
\toprule
\textbf{Knoten} & \textbf{Ab Rück} & \textbf{An Rück} & \textbf{Berechnung} \\
\midrule
G & 10 & -- & Startzeit \\
F & 13 & $(10+2)=12$ & \\
E & 17 & $(13+3)=16$ & \\
D & 25 & $(17+7)=24$ & \\
C & 32 & $(25+6)=31$ & Abschnitt C$\to$B: Rück fährt ab 32 \\
B & 41 & $(32+8)=40$ & Rück-Intervall C$\to$B: $[32, 40]$ \\
A & -- & $(41+5)=46$ & \\
\bottomrule
\end{tabular}
\end{center}

Kollisionsprüfung B$\to$C: Hin $[23,31]$, Rück $[32, 40]$ -- \textbf{kein Überlapp} $\Rightarrow$ keine Kollision! Strafe = 0.

Da 0 die minimale mögliche Strafe ist, wird \texttt{besteStartzeitRueck = 10} gespeichert.

\textbf{Schritt 2: Warteverteilung}

Da alle \texttt{wartezeitRueck = 0} sind, ändert \texttt{beidseitigesWarten()} nichts.

\textbf{Endergebnis}:
\begin{center}
\begin{tabular}{r|ccccccc}
 & A & B & C & D & E & F & G \\
\hline
An Hin  &    & 22 & 31 & 38 & 46 & 50 & 53 \\
Wa Hin  &    &    &    &    &    &    &    \\
Ab Hin  & 17 & 23 & 32 & 39 & 47 & 51 &    \\
\hline
Bahnhof & A  & B  & C  & D  & E  & F  & G  \\
\hline
Ab Rück &    & 41 & 32 & 25 & 17 & 13 & 10 \\
Wa Rück &    &    &    &    &    &    &    \\
An Rück & 46 & 40 & 31 & 24 & 16 & 12 &    \\
\end{tabular}
\end{center}

Gesamtdauer Hin/Rück: 36 / 36 \quad Wartezeit: 0 / 0 \quad Strafe: $\mathbf{0}$
\end{beispielbox}

\subsection{Berechnung: Beispiel3 (Fall mit Wartezeit-Aufteilung)}

\begin{beispielbox}[title={Beidseitiges Warten -- Beispiel3: A--B--C--D--E--F--G--H--I--J, Start: 17}]

Beispiel3 hat zwei Kollisionsstellen (A und D). Die Suchschleife findet eine optimale Startzeit, bei der noch Restkollisionen verbleiben, die dann hälftig aufgeteilt werden.

Tatsächliche Ausgabe des Programms:
\begin{center}
\begin{tabular}{r|cccccccccc}
 & A & B & C & D & E & F & G & H & I & J \\
\hline
An Hin  &   & 29 & 38 & 53 & 57 & 03 & 25 & 39 & 46 & 55 \\
Wa Hin  & 02 &    &    & 03 &    &    &    &    &    &   \\
Ab Hin  & 17 & 30 & 39 & 54 & 01 & 08 & 30 & 44 & 51 & 00 \\
\hline
Bahn.   & A  & B  & C  & D  & E  & F  & G  & H  & I  & J \\
\hline
Ab Rück &    & 30 & 19 & 04 & 00 & 52 & 30 & 16 & 09 & 60 \\
Wa Rück & 02 &    &    & 02 &    &    &    &    &    &   \\
An Rück & 42 & 27 & 18 & 03 & 57 & 51 & 29 & 15 & 08 &   \\
\end{tabular}
\end{center}

Gesamtdauer Hin/Rück: 102 / 102 \quad Wartezeit: 4 / 4 \quad Strafe: $4^2 + 4^2 = \mathbf{32}$

\textbf{Aufteilung im Vergleich}:

\begin{center}
\begin{tabular}{lccc}
\toprule
 & \textbf{Wartezeit Hin} & \textbf{Wartezeit Rück} & \textbf{Strafe} \\
\midrule
Einseitiges Warten & 0 & 8 & $8^2 = 64$ \\
Beidseitiges Warten & 4 & 4 & $4^2 + 4^2 = 32$ \\
\bottomrule
\end{tabular}
\end{center}

Die Gesamtwartezeit ist in beiden Fällen 8 Minuten. Durch die gleichmäßige Aufteilung halbiert sich die Strafe von 64 auf 32 -- exakt wie die Theorie vorhersagt: $\frac{W^2}{2} = \frac{64}{2} = 32$.
\end{beispielbox}

\subsection{Vergleich der drei Strategien an Beispiel2}

\begin{center}
\begin{tabular}{lccc}
\toprule
\textbf{Strategie} & \textbf{Gesamtdauer Hin/Rück} & \textbf{Gesamtwartezeit} & \textbf{Strafe} \\
\midrule
Einfache Fahrt & 36 / 36 & 0 / 0 & 0 \\
Einseitiges Warten & 36 / 52 & 0 / 16 & 256 \\
Beidseitiges Warten & 36 / 36 & 0 / 0 & 0 \\
\bottomrule
\end{tabular}
\end{center}

\begin{hinweisbox}
Die Einfache Fahrt zeigt Strafe 0, weil sie Kollisionen gar nicht behebt und keine Wartezeiten einbaut -- sie ignoriert das Problem. Das Beidseitige Warten löst die Kollision korrekt auf, indem es die Rückfahrt so startet, dass sich die Züge nicht begegnen.
\end{hinweisbox}

% ─────────────────────────────────────────
\section{Ausgabe (AusgabeManager)}

\subsection{Überblick}

Die Klasse \texttt{AusgabeManager} übernimmt die gesamte Formatierung und Speicherung der Ergebnisse. Sie enthält vier Methoden:

\begin{center}
\begin{tabular}{ll}
\toprule
\textbf{Methode} & \textbf{Aufgabe} \\
\midrule
\texttt{druckeErgebnis(strategie, plan)} & Baut den Ausgabe-String auf und gibt ihn zurück \\
\texttt{printRow(label, extraktor, plan, isStation)} & Erzeugt eine einzelne Zeile der Tabelle \\
\texttt{berechneMindestdauer(plan)} & Berechnet die Fahrzeit ohne Wartezeiten \\
\texttt{speicherErgebnis(dateiname, inhalt)} & Schreibt den Ausgabe-String in eine Datei \\
\bottomrule
\end{tabular}
\end{center}

\subsection{druckeErgebnis()}

Diese Methode baut den gesamten Ausgabe-String für eine Strategie zusammen:

\begin{lstlisting}[caption={AusgabeManager.druckeErgebnis()}]
public static String druckeErgebnis(FahrplanStrategie strategie, BahnhofPlan plan) {
    StringBuilder sb = new StringBuilder();
    sb.append(strategie.getStrategieName()).append(":\n");

    // Hinfahrt-Zeilen
    sb.append(printRow("An ", n -> n.getAnkunfthin() != null
        ? String.format("%02d", n.getAnkunfthin()) : "  ", plan, false));
    sb.append(printRow("Wa ", n -> n.getWartezeitHin() > 0
        ? String.format("%02d", n.getWartezeitHin()) : "  ", plan, false));
    sb.append(printRow("Ab ", n -> n.getAbfahrthin() != null
        ? String.format("%02d", n.getAbfahrthin()) : "  ", plan, false));

    // Bahnhofnamen-Zeile (mit Kollisionsmarkierung)
    sb.append(printRow("   ", n -> String.format("%2s", n.getName()), plan, true));

    // Rueckfahrt-Zeilen
    sb.append(printRow("Ab ", n -> n.getAbfahrtrueck() != null
        ? String.format("%02d", n.getAbfahrtrueck()) : "  ", plan, false));
    sb.append(printRow("Wa ", n -> n.getWartezeitRueck() > 0
        ? String.format("%02d", n.getWartezeitRueck()) : "  ", plan, false));
    sb.append(printRow("An ", n -> n.getAnkunftrueck() != null
        ? String.format("%02d", n.getAnkunftrueck()) : "  ", plan, false));

    // Statistik
    int mindestdauer = berechneMindestdauer(plan);
    int warteHin = 0, warteRueck = 0;
    BahnhofNode current = plan.getHead();
    while (current != null) {
        warteHin  += current.getWartezeitHin();
        warteRueck += current.getWartezeitRueck();
        current = current.getNext();
    }
    sb.append(String.format("Gesamtdauer Hinfahrt, Rueckfahrt       : %d, %d%n",
        mindestdauer + warteHin, mindestdauer + warteRueck));
    sb.append(String.format("Summe Wartezeiten Hinfahrt, Rueckfahrt : %d, %d%n",
        warteHin, warteRueck));
    sb.append(String.format("Summe Strafen                         : %d%n",
        strategie.getSummeStrafen()));
    return sb.toString();
}
\end{lstlisting}

Die Methode arbeitet rein auf Lesezugriff -- sie verändert keine Daten im \texttt{BahnhofPlan}.

\subsection{printRow() -- Hilfsmethode für horizontale Zeilen}

Das besondere Layout der Ausgabe (Bahnhöfe als Spalten, Zeiten als Zeilen) wird durch \texttt{printRow()} realisiert:

\begin{lstlisting}[caption={AusgabeManager.printRow() mit Funktionsinterface}]
private static String printRow(
        String label,
        Function<BahnhofNode, String> wertExtraktor,
        BahnhofPlan plan,
        boolean isStationRow) {
    StringBuilder row = new StringBuilder();
    row.append(String.format("%-3s", label));
    BahnhofNode current = plan.getHead();
    while (current != null) {
        row.append(wertExtraktor.apply(current));
        if (current.getNext() != null) {
            // 'x' zwischen zwei Bahnhofnamen bei Kollision
            row.append(isStationRow && current.isKollision() ? "x " : "  ");
        }
        current = current.getNext();
    }
    row.append("\n");
    return row.toString();
}
\end{lstlisting}

Der Schlüssel liegt im Parameter \texttt{wertExtraktor}: Er ist ein \texttt{Function<BahnhofNode, String>} -- also ein Java-Lambda, das für jeden Knoten den anzuzeigenden String bestimmt. Dadurch kann \texttt{printRow()} für \emph{alle} sieben Zeilen (An-Hin, Wa-Hin, Ab-Hin, Bahnhöfe, Ab-Rück, Wa-Rück, An-Rück) wiederverwendet werden, ohne die Schleife über die Liste zu duplizieren.

Beispiel für den Aufruf der Ankunftszeile der Hinfahrt:
\begin{lstlisting}
printRow("An ", n -> n.getAnkunfthin() != null
    ? String.format("%02d", n.getAnkunfthin()) : "  ", plan, false);
\end{lstlisting}

Das Lambda gibt \texttt{"  "} (zwei Leerzeichen) zurück, wenn \texttt{ankunfthin == null} ist (z.\,B. beim ersten Bahnhof, der keine Ankunftszeit hat), andernfalls die zweistellige Minutenangabe.

\subsection{berechneMindestdauer()}

\begin{lstlisting}[caption={Berechnung der Mindestfahrtdauer ohne Wartezeiten}]
public static int berechneMindestdauer(BahnhofPlan plan) {
    int dauer = 0;
    int anzahlBahnhoefe = 0;
    BahnhofNode current = plan.getHead();
    while (current != null) {
        anzahlBahnhoefe++;
        if (current.getNext() != null)
            dauer += current.getDistanceToNext();
        current = current.getNext();
    }
    // 1 Min Haltezeit fuer jeden Zwischenbahnhof
    if (anzahlBahnhoefe > 2)
        dauer += (anzahlBahnhoefe - 2);
    return dauer;
}
\end{lstlisting}

Die Mindestdauer ist die Summe aller Streckenabstände plus je eine Minute Haltezeit pro Zwischenbahnhof (Start- und Endbahnhof zählen nicht). Sie wird von jeder Strategie als gemeinsame Basis genutzt; die tatsächliche Gesamtdauer ergibt sich als $\texttt{Mindestdauer} + \texttt{SummeWartezeiten}$.

\subsection{speicherErgebnis()}

\begin{lstlisting}[caption={Ergebnis in Datei schreiben}]
public static void speicherErgebnis(String dateiname, String inhalt) {
    File verzeichnis = new File("ausgabe");
    if (!verzeichnis.exists()) verzeichnis.mkdir();
    File zielDatei = new File(verzeichnis,
        dateiname.replace(".txt", "") + "_ergebnis.txt");
    try (BufferedWriter writer = new BufferedWriter(new FileWriter(zielDatei))) {
        writer.write(inhalt);
    } catch (IOException e) {
        System.err.println("Fehler beim Speichern: " + e.getMessage());
    }
}
\end{lstlisting}

Das Verzeichnis \texttt{ausgabe/} wird beim ersten Aufruf automatisch angelegt, falls es noch nicht existiert. Der Dateiname ergibt sich aus dem Eingabedateinamen mit dem Suffix \texttt{\_ergebnis.txt}.

\subsection{Format der Ausgabe}

Die erzeugte Ausgabe für Beispiel2, Strategie ``Einseitiges Warten'', sieht folgendermaßen aus:

\begin{verbatim}
Einseitiges Warten:
An     22  31  38  46  50  53
Wa
Ab 17  23  32  39  47  51
    A   B   C   D   E   F   G
Ab     41  32  09  01  57  54
Wa         16
An 46  40  15  08  00  56
Gesamtdauer Hinfahrt, Rueckfahrt       : 36, 52
Summe Wartezeiten Hinfahrt, Rueckfahrt : 0, 16
Summe Strafen                          : 256
\end{verbatim}

\begin{itemize}
    \item \textbf{Zeile 1}: Strategiename
    \item \textbf{An / Wa / Ab (oben)}: Ankunft-, Wartezeit- und Abfahrtszeile der Hinfahrt
    \item \textbf{Bahnhofzeile}: Namen mit \texttt{x} für Kollisionsabschnitte
    \item \textbf{Ab / Wa / An (unten)}: Abfahrt-, Wartezeit- und Ankunftszeile der Rückfahrt
    \item \textbf{Statistikzeilen}: Gesamtdauer, Wartezeiten und Strafe
\end{itemize}

% ─────────────────────────────────────────
\section{UML-Diagramme (Mermaid)}

Die folgenden Abschnitte enthalten die Mermaid-Quellcodes der UML-Diagramme. Sie können unter \texttt{https://mermaid.live} oder in kompatiblen Markdown-Editoren gerendert werden.

\subsection{Klassendiagramm}

\begin{lstlisting}[style=mermaid, caption={Klassendiagramm -- vollständige Klassenstruktur}]
classDiagram
    class BahnhofNode {
        -String name
        -Integer ankunfthin
        -Integer abfahrthin
        -Integer wartezeitHin
        -Integer ankunftrueck
        -Integer abfahrtrueck
        -Integer wartezeitRueck
        -BahnhofNode next
        -BahnhofNode prev
        -int distanceToNext
        +isKollision() boolean
        -checkOverlap(s1,e1,s2,e2) boolean
        -isTimeInInterval(t,start,ende) boolean
    }

    class BahnhofPlan {
        -BahnhofNode head
        -BahnhofNode tail
        -int groesse
        +addBahnhof(name) void
        +setAbstaende(abstaende) void
        +reset() void
        +BahnhofPlan(original) -- Kopierkonstruktor
    }

    class SimulationsDaten {
        <<record>>
        +BahnhofPlan plan
        +int startZeit
    }

    class FahrplanStrategie {
        <<abstract>>
        #BahnhofPlan aktuellerPlan
        +getStrategieName()* String
        +berechneFahrplan(plan, startZeit)* void
        +getSummeStrafen()* int
        +getSummeWarteZeit()* int
        +hinfahrtOhneKollision(plan, startZeit)$ void
        +rueckfahrtOhneKollision(plan, startZeit)$ void
        +hinfahrtMitKollision(plan, startZeit)$ void
        +rueckfahrtMitKollision(plan, startZeit)$ void
        +beidseitigesWarten(plan, startZeit)$ void
        +warteZeitHinReset(plan)$ void
        +warteZeitRueckReset(plan)$ void
    }

    class EinfacheFahrtStrategie {
        +berechneFahrplan(plan, startZeit) void
        +getSummeStrafen() int
        +getSummeWarteZeit() int
    }

    class EinseitigesWartenStrategie {
        +berechneFahrplan(plan, startZeit) void
        +getSummeStrafen() int
        +getSummeWarteZeit() int
    }

    class BeidseitigesWartenStrategie {
        -int besteStrafe
        -int besteStartzeitRueck
        +berechneFahrplan(plan, startZeit) void
        +getSummeStrafen() int
        +getSummeWarteZeit() int
    }

    class EingabeLeser {
        +leseDatei(dateiPfad)$ SimulationsDaten
    }

    class AusgabeManager {
        +druckeErgebnis(strategie, plan)$ String
        +berechneMindestdauer(plan)$ int
        +speicherErgebnis(dateiname, inhalt)$ void
        -printRow(label, extraktor, plan, isStation)$ String
    }

    FahrplanStrategie <|-- EinfacheFahrtStrategie
    FahrplanStrategie <|-- EinseitigesWartenStrategie
    FahrplanStrategie <|-- BeidseitigesWartenStrategie
    BahnhofPlan "1" *-- "1..*" BahnhofNode : enthaelt
    BahnhofNode --> BahnhofNode : next / prev
    SimulationsDaten --> BahnhofPlan
    EingabeLeser ..> SimulationsDaten : erstellt
    AusgabeManager ..> BahnhofPlan : liest
    AusgabeManager ..> FahrplanStrategie : liest
\end{lstlisting}

\subsection{Sequenzdiagramm: Hauptprogrammablauf}

\begin{lstlisting}[style=mermaid, caption={Sequenzdiagramm -- Ablauf in Main.java}]
sequenceDiagram
    participant Main
    participant EingabeLeser
    participant BahnhofPlan
    participant S1 as EinfacheFahrtStrategie
    participant S2 as EinseitigesWartenStrategie
    participant S3 as BeidseitigesWartenStrategie
    participant AusgabeManager

    Main->>EingabeLeser: leseDatei(pfad)
    EingabeLeser->>BahnhofPlan: addBahnhof() x n
    EingabeLeser->>BahnhofPlan: setAbstaende()
    EingabeLeser-->>Main: SimulationsDaten(plan, startZeit)

    Note over Main: Gleicher BahnhofPlan fuer alle Strategien!

    Main->>S1: berechneFahrplan(plan, startZeit)
    S1->>BahnhofPlan: schreibt ankunfthin, abfahrthin (Hinfahrt)
    S1->>BahnhofPlan: schreibt ankunftrueck, abfahrtrueck (Rueckfahrt)
    Main->>AusgabeManager: druckeErgebnis(S1, plan)
    AusgabeManager-->>Main: String

    Main->>S2: berechneFahrplan(plan, startZeit)
    Note over S2: Liest ankunfthin(tail) von S1!
    S2->>BahnhofPlan: ueberschreibt Rueckfahrtdaten + wartezeitRueck
    Main->>AusgabeManager: druckeErgebnis(S2, plan)
    AusgabeManager-->>Main: String

    Main->>S3: berechneFahrplan(plan, startZeit)
    Note over S3: Liest ankunfthin(tail) von S1!
    loop alle 60 Startzeiten
        S3->>BahnhofPlan: rueckfahrtMitKollision(plan, i)
        S3->>S3: Strafe berechnen und vergleichen
    end
    S3->>BahnhofPlan: beste Ruckfahrt + beidseitigesWarten
    Main->>AusgabeManager: druckeErgebnis(S3, plan)
    AusgabeManager-->>Main: String

    Main->>AusgabeManager: speicherErgebnis(dateiname, gesamtausgabe)
\end{lstlisting}

\subsection{Ablaufdiagramm: rueckfahrtMitKollision()}

\begin{lstlisting}[style=mermaid, caption={Flussdiagramm -- Algorithmus rueckfahrtMitKollision()}]
flowchart TD
    A([Start: rueckfahrtMitKollision\nplan, startZeit]) --> B

    B[Phase 1: rueckfahrtOhneKollision\nBerechnet Ausgangsbasis ohne Kollision] --> C

    C[current = tail.prev\nStarte Rueckwaertstraversierung] --> D

    D{current.prev\n!= null?} -- Nein --> Z

    D -- Ja --> E

    E["current.ankunftrueck =
    (current.next.abfahrtrueck
    + current.distanceToNext) % 60"]

    E --> F["current.abfahrtrueck =
    (current.ankunftrueck + 1) % 60"]

    F --> G["current.prev.ankunftrueck =
    (current.abfahrtrueck
    + current.prev.distanceToNext) % 60"]

    G --> H{current.prev\n.isKollision?}

    H -- Nein --> J[wartezeitRueck bleibt 0]

    H -- Ja --> I["wartezeitRueck =
    (ankunfthin - ankunftrueck + 60) % 60"]

    I --> K

    J --> K["abfahrtrueck =
    (ankunftrueck + wartezeitRueck + 1) % 60"]

    K --> L[current = current.prev]
    L --> D

    Z["A.ankunftrueck =
    (B.abfahrtrueck + A.distanceToNext) % 60"]
    Z --> END([Ende])
\end{lstlisting}

\subsection{Ablaufdiagramm: BeidseitigesWartenStrategie}

\begin{lstlisting}[style=mermaid, caption={Flussdiagramm -- Algorithmus BeidseitigesWartenStrategie.berechneFahrplan()}]
flowchart TD
    A([Start: berechneFahrplan\nplan, startZeitHin]) --> B

    B["besteStrafe = MAX_INT\nbesteStartzeitRueck = -1\ni = 0"] --> C

    C{i < 60?} -- Nein --> H

    C -- Ja --> D["rueckfahrtMitKollision(plan, i)"]

    D --> E{tail.prev\nhat Kollision?}

    E -- Ja: ungueltig --> F[i++]
    F --> C

    E -- Nein --> G["strafe = getSummeStrafen()"]

    G --> G2{strafe <\nbesteStrafe?}

    G2 -- Nein --> F

    G2 -- Ja --> G3["besteStrafe = strafe\nbesteStartzeitRueck = i"]
    G3 --> F

    H["rueckfahrtMitKollision\n(plan, besteStartzeitRueck)"] --> I

    I{besteStartzeitRueck\n!= 0?} -- Nein --> END

    I -- Ja --> J["beidseitigesWarten(plan, startZeitHin)"]

    J --> END([Ende])
\end{lstlisting}

\subsection{Ablaufdiagramm: Kollisionserkennung (isKollision)}

\begin{lstlisting}[style=mermaid, caption={Flussdiagramm -- BahnhofNode.isKollision()}]
flowchart TD
    A([Start: isKollision\nan Knoten B]) --> B

    B{next == null?} -- Ja --> NEIN([return false])

    B -- Nein --> C

    C["hStart = this.abfahrthin\nhEnde  = next.ankunfthin\nrStart = next.abfahrtrueck\nrEnde  = this.ankunftrueck"]

    C --> D["checkOverlap(hStart, hEnde, rStart, rEnde)"]

    D --> E[m = 0]

    E --> F{m < 60?} -- Nein --> NEIN2([return false])

    F -- Ja --> G{isTimeInInterval\nm in Hin-Intervall?}

    G -- Nein --> L[m++]

    G -- Ja --> H{isTimeInInterval\nm in Rueck-Intervall?}

    H -- Nein --> L

    H -- Ja --> JA([return true\nKOLLISION!])

    L --> F
\end{lstlisting}

\subsection{Ablaufdiagramm: EingabeLeser}

\begin{lstlisting}[style=mermaid, caption={Flussdiagramm -- EingabeLeser.leseDatei()}]
flowchart TD
    A([Start: leseDatei\ndateiPfad]) --> B

    B["zeilen = Files.readAllLines(pfad)\nbahnhofPlan = new BahnhofPlan()\nstartZeit = 0\ni = 0"]

    B --> C{i < zeilen\n.size?} -- Nein --> Z

    C -- Ja --> D["aktuelleZeile = zeilen.get(i).trim()"]

    D --> E{Strecke:?}

    E -- Ja --> F["zeilen[i+1] splitten\nfuer jeden Namen:\nbahnhofPlan.addBahnhof(name)"]
    F --> L[i++]

    E -- Nein --> G{Abstaende:?}

    G -- Ja --> H["zeilen[i+1] splitten\nals int[] parsen\nbahnhofPlan.setAbstaende()"]
    H --> L

    G -- Nein --> I{"Start\nHinfahrt:?"}

    I -- Ja --> J["startZeit =\nInteger.parseInt(zeilen[i+1])"]
    J --> L

    I -- Nein --> L
    L --> C

    Z["return new SimulationsDaten\n(bahnhofPlan, startZeit)"]
    Z --> END([Ende])
\end{lstlisting}

% ─────────────────────────────────────────
\section{Testfälle: Normalfälle, Grenzfalltests und Fehlertests}
\label{sec:tests}

Dieses Kapitel dokumentiert alle Testfälle des Projekts. Jeder Testfall enthält
Eingabe, erwartete Ausgabe (Ausschnitt), Analyse und — bei Fehlertests — den
beobachteten Fehler.

\medskip
\noindent\textbf{Kategorien:}
\begin{itemize}
  \item \textbf{N} – Normalfall: typische Eingaben ohne Sonderbedingungen
  \item \textbf{G} – Grenzfall: Extreme Eingabewerte (Min/Max)
  \item \textbf{F} – Fehlerfall: Fehlerhafte Eingaben oder bekannte Programmfehler
  \item \textbf{S} – Spezialfall: Besondere Szenarien (Kaskadeneffekte, Mitternacht)
\end{itemize}

% ══════════════════════════════════════════════════════════════════════════════
% SEITE 1 – N1 und N2
% ══════════════════════════════════════════════════════════════════════════════

\begin{tcolorbox}[
  title={\textbf{N1 – Basisfall: Keine Kollision} \hfill\texttt{01\_beispiel1\_baseline.txt}},
  colback=green!4!white, colframe=green!50!black, fonttitle=\bfseries,
  breakable, before upper={\parindent0pt}]

\textbf{Eingabe:}
\begin{verbatim}
Strecke:    A B C
Abstaende:  5 7
Start:      17
\end{verbatim}

\textbf{Analyse:}\\
Hinfahrt: $A \xrightarrow{5} B \xrightarrow{7} C$. Abfahrt A=17, Ankunft B=22,
Abfahrt B=23, Ankunft C=30, Abfahrt C=31. Rückfahrt startet bei C=31.
A→B und B→A fahren zu unterschiedlichen Zeiten (kein Überlapp) → keine Kollision.
Alle drei Strategien liefern dasselbe Ergebnis; Strafe = 0.

\textbf{Ausgabe (Auszug):}
\begin{verbatim}
An     22  30
Ab 17  23  31
    A   B   C
Ab 45  39  31
An 44  38
Strafe: 0
\end{verbatim}

\textbf{Erwartetes Ergebnis:} Alle Wartezeiten = 0, keine Kollisionsmarkierung.\\
\textbf{Tatsächliches Ergebnis:} \textcolor{green!60!black}{\checkmark\ Korrekt.}

\end{tcolorbox}

\begin{tcolorbox}[
  title={\textbf{N2 – Eine Kollision} \hfill\texttt{02\_beispiel2\_eineKollision.txt}},
  colback=blue!4!white, colframe=blue!50!black, fonttitle=\bfseries,
  breakable, before upper={\parindent0pt}]

\textbf{Eingabe:}
\begin{verbatim}
Strecke:    A B C D E F G
Abstaende:  5 9 6 7 3 2
Start:      17
\end{verbatim}

\textbf{Analyse:}\\
Segment A→B: Hin fährt [17,22], Rück fährt [27,33].
Segment B→C: Hin [23,32], Rück [17,27] → Überlapp bei Minute 23–27 → \textbf{Kollision bei B}.

\begin{itemize}
  \item \textbf{Einfache Fahrt:} Kollision markiert (\texttt{Bx}), Strafe = 0.
  \item \textbf{Einseitiges Warten:} C wartet 16 min (Rück), Strafe = $16^2 = 256$.
  \item \textbf{Beidseitiges Warten:} Rückfahrt startet bei 49, keine Kollision, Strafe = 0.
\end{itemize}

\textbf{Ausgabe (Einseitiges Warten):}
\begin{verbatim}
    A   B   C   D   E   F   G
Ab 33  43  33  10  02  58  55
Wa         16
Summe Strafen: 256
\end{verbatim}

\textbf{Erwartetes Ergebnis:} Einseitiges Warten = 256, Beidseitiges = 0.\\
\textbf{Tatsächliches Ergebnis:} \textcolor{green!60!black}{\checkmark\ Korrekt.}

\end{tcolorbox}

\newpage
% ══════════════════════════════════════════════════════════════════════════════
% SEITE 2 – N3 und G1
% ══════════════════════════════════════════════════════════════════════════════

\begin{tcolorbox}[
  title={\textbf{N3 – Zehn Stationen} \hfill\texttt{03\_beispiel3\_zehnStationen.txt}},
  colback=green!4!white, colframe=green!50!black, fonttitle=\bfseries,
  breakable, before upper={\parindent0pt}]

\textbf{Eingabe:}
\begin{verbatim}
Strecke:    A B C D E F G H I J
Abstaende:  12 8 14 3 5 21 13 6 8
Start:      17
\end{verbatim}

\textbf{Analyse:}\\
Zehn Stationen, neun Abstände (Summe = 90). Die Hinfahrt dauert
$90 + 8 = 98$ Minuten (inkl. 8 Haltezeiten). Auf einer Modulo-60-Uhr entspricht
das 38 Minuten nach Start. Durch die vielen Überlappungsbereiche prüft das
Programm alle Segmente auf Kollision.

Keine Kollision erwartet (unterschiedliche Abstände erzeugen zeitlich getrennte
Intervalle). Alle drei Strategien liefern identische Ergebnisse, Strafe = 0.

\textbf{Erwartetes Ergebnis:} Keine Kollision, Strafe = 0 bei allen Strategien.\\
\textbf{Tatsächliches Ergebnis:} \textcolor{green!60!black}{\checkmark\ Korrekt.}

\end{tcolorbox}

\begin{tcolorbox}[
  title={\textbf{G1 – Minimale Stationsanzahl: 2 Stationen} \hfill\texttt{04\_grenz\_zweiStationen.txt}},
  colback=orange!5!white, colframe=orange!60!black, fonttitle=\bfseries,
  breakable, before upper={\parindent0pt}]

\textbf{Eingabe:}
\begin{verbatim}
Strecke:    A B
Abstaende:  29
Start:      0
\end{verbatim}

\textbf{Analyse:}\\
Zwei Stationen sind die kleinstmögliche Eingabe – eine einzige Strecke A→B.
Abstand = 29: Hin fährt [0,29], Rück startet bei 30 und fährt [30,59].
Die Intervalle berühren sich exakt bei Minute 29/30. Da \texttt{isKollision()} nur
Überlapp (beide zur selben Minute auf dem Abschnitt) prüft, ist kein Überlapp
vorhanden: Hin endet bei 29 (Ankunft B), Rück startet bei 30 (Abfahrt B).

\textbf{Ausgabe:}
\begin{verbatim}
An     29         An     29         An     29
Ab 00  30         Ab 00  30         Ab 00  30
    A   B             A   B             A   B
Ab 00  30         Ab 00  30         Ab 00  30
An 59             An 59             An 59
Strafe: 0         Strafe: 0         Strafe: 0
\end{verbatim}

\textbf{Erwartetes Ergebnis:} Kein Fehler trotz minimalem Graphen, Strafe = 0.\\
\textbf{Tatsächliches Ergebnis:} \textcolor{green!60!black}{\checkmark\ Korrekt.}

\end{tcolorbox}

\newpage
% ══════════════════════════════════════════════════════════════════════════════
% SEITE 3 – G2 und G3
% ══════════════════════════════════════════════════════════════════════════════

\begin{tcolorbox}[
  title={\textbf{G2 – Startzeit = 59 (Modulo-Wrap-around)} \hfill\texttt{05\_grenz\_startzeit59\_wraparound.txt}},
  colback=orange!5!white, colframe=orange!60!black, fonttitle=\bfseries,
  breakable, before upper={\parindent0pt}]

\textbf{Eingabe:}
\begin{verbatim}
Strecke:    A B C
Abstaende:  5 7
Start:      59
\end{verbatim}

\textbf{Analyse:}\\
Startzeit 59 ist der höchste gültige Wert der Modulo-60-Uhr. Die Abfahrt von A
erfolgt bei Minute 59. Die nächste Abfahrt wäre Minute 64 — auf der
Modulo-60-Uhr ist das Minute 4. Die Berechnungen:

\begin{center}
\begin{tabular}{lll}
\toprule
Bahnhof & Ankunft & Abfahrt \\
\midrule
A       & —       & 59      \\
B       & $(59+5)\%60 = 4$ & 5 \\
C       & $(5+7)\%60 = 12$ & 13 \\
\bottomrule
\end{tabular}
\end{center}

Der Wrap-around (59 → 4) wird korrekt durch Modulo-60-Arithmetik berechnet.
Keine Kollision, Strafe = 0.

\textbf{Erwartetes Ergebnis:} Wrap-around korrekt behandelt, Strafe = 0.\\
\textbf{Tatsächliches Ergebnis:} \textcolor{green!60!black}{\checkmark\ Korrekt.}

\end{tcolorbox}

\begin{tcolorbox}[
  title={\textbf{G3 – Startzeit = 0 (untere Grenze)} \hfill\texttt{09\_grenz\_startzeit0.txt}},
  colback=orange!5!white, colframe=orange!60!black, fonttitle=\bfseries,
  breakable, before upper={\parindent0pt}]

\textbf{Eingabe:}
\begin{verbatim}
Strecke:    A B C D
Abstaende:  15 15 15
Start:      0
\end{verbatim}

\textbf{Analyse:}\\
Startzeit 0 ist die kleinstmögliche gültige Startzeit. Segment A→B: Hin [0,15],
Rück startet bei 48 und fährt [48,3] (Wrap-around). B→C: Hin [16,31], Rück [32,48].
\textbf{Überlapp bei B→C}: Hin endet 31, Rück startet 32 → kein Überlapp?
Nein: Hin [16,31], Rück [48,3 wraparound]. Kollision bei A→B:
Rück [48,3] überschneidet Hin [0,15] → Kollision.

Einseitiges Warten erzeugt große Wartezeit (56 min), Strafe = $56^2 = 3136$.
Beidseitiges Warten teilt auf: je 14 min, Strafe = $14^2 + 14^2 = 392$.

\textbf{Ausgabe (Beidseitiges Warten):}
\begin{verbatim}
Wa     14
Ab 00  30  46  02
    A   B   C   D
Ab 18  16  46  30
Wa     14
Strafe: 392
\end{verbatim}

\textbf{Erwartetes Ergebnis:} Strafe einseitig 3136, beidseitig 392.\\
\textbf{Tatsächliches Ergebnis:} \textcolor{green!60!black}{\checkmark\ Korrekt.}

\end{tcolorbox}

\newpage
% ══════════════════════════════════════════════════════════════════════════════
% SEITE 4 – G4 und G5
% ══════════════════════════════════════════════════════════════════════════════

\begin{tcolorbox}[
  title={\textbf{G4 – Minimaler Abstand = 1} \hfill\texttt{12\_minimalAbstand1.txt}},
  colback=orange!5!white, colframe=orange!60!black, fonttitle=\bfseries,
  breakable, before upper={\parindent0pt}]

\textbf{Eingabe:}
\begin{verbatim}
Strecke:    A B C D E F
Abstaende:  1 1 1 1 1
Start:      0
\end{verbatim}

\textbf{Analyse:}\\
Alle Abstände = 1 ist der kleinstmögliche Abstandswert. Hinfahrt dauert
$5 \times 1 + 4 = 9$ Minuten (5 Segmente, 4 Haltezeiten). Rückfahrt startet bei
Minute 10. Hin und Rück sind zeitlich weit getrennt (Rück startet bei 10, Hin
endet bei 9) → keine Kollision.

\textbf{Ausgabe:}
\begin{verbatim}
An     01  03  05  07  09
Ab 00  02  04  06  08  10
    A   B   C   D   E   F
Ab 20  18  16  14  12  10
An 19  17  15  13  11
Strafe: 0 (alle Strategien)
\end{verbatim}

\textbf{Erwartetes Ergebnis:} Kein Fehler bei Abstand = 1, Strafe = 0.\\
\textbf{Tatsächliches Ergebnis:} \textcolor{green!60!black}{\checkmark\ Korrekt.}

\end{tcolorbox}

\begin{tcolorbox}[
  title={\textbf{G5 – Stresstest: 52 Stationen} \hfill\texttt{13\_stress\_test.txt}},
  colback=orange!5!white, colframe=orange!60!black, fonttitle=\bfseries,
  breakable, before upper={\parindent0pt}]

\textbf{Eingabe:}
\begin{verbatim}
Strecke:    A B C ... Z AA AB ... AZ  (52 Stationen)
Abstaende:  4 11 7 18 3 22 9 ...     (51 Abstände)
Start:      37
\end{verbatim}

\textbf{Analyse:}\\
Dieser Test prüft zwei Eigenschaften gleichzeitig:
\begin{enumerate}
  \item \textbf{Zweistellige Stationsnamen} (\texttt{AA}–\texttt{AZ}): Erfordert
        \texttt{String} statt \texttt{char} für den Namen (Grenzfall der
        Entwurfsänderung aus Kap.~3.1).
  \item \textbf{Laufzeit- und Speicherverhalten} bei $n = 52$: Alle Algorithmen
        müssen auch mit vielen Knoten korrekt traversieren.
\end{enumerate}

Erwartet wird, dass das Programm terminiert, die Ausgabe alle 52 Stationen korrekt
tabellarisch ausgibt und keine \texttt{NullPointerException} beim Traversieren
des Endes der Liste auftritt.

\textbf{Erwartetes Ergebnis:} Korrekte Terminierung, keine Exception,
zweistellige Namen in der Ausgabe.\\
\textbf{Tatsächliches Ergebnis:} \textcolor{green!60!black}{\checkmark\ Korrekt.}

\end{tcolorbox}

\newpage
% ══════════════════════════════════════════════════════════════════════════════
% SEITE 5 – S1 und S2
% ══════════════════════════════════════════════════════════════════════════════

\begin{tcolorbox}[
  title={\textbf{S1 – Mitternacht-Überschreitung (Rollover)} \hfill\texttt{11\_rollover\_mitternacht.txt}},
  colback=violet!5!white, colframe=violet!50!black, fonttitle=\bfseries,
  breakable, before upper={\parindent0pt}]

\textbf{Eingabe:}
\begin{verbatim}
Strecke:    A B C
Abstaende:  25 25
Start:      55
\end{verbatim}

\textbf{Analyse:}\\
Segment A→B: Hin fährt [55, 20] (überschreitet Minute 59 → 0). Das Intervall
umschließt den Tageswechsel.

\begin{center}
\begin{tabular}{lll}
\toprule
Bahnhof & Ankunft & Abfahrt \\
\midrule
A       & —       & 55      \\
B       & $(55+25)\%60 = 20$ & 21 \\
C       & $(21+25)\%60 = 46$ & 47 \\
\bottomrule
\end{tabular}
\end{center}

Rück: Abfahrt C=47, Ankunft B=22, Abfahrt B=21, Ankunft A=55 (wraparound).\\
Segment A→B: Hin [55,20], Rück [21,47]. Hin umschließt [55..59, 0..20],
Rück läuft [21..47]. Überlapp? Hin endet bei 20, Rück startet bei 21 → kein Überlapp → \textbf{Kollision!}
Minutengenaue Prüfung: bei Minute 20 fährt Hin noch (kommt an), Rück ist noch
auf Abfahrt → Kollision erkannt an Segment A (markiert \texttt{Ax}).

\textbf{Ausgabe (Einseitiges Warten):}
\begin{verbatim}
Ab 55  21  47
    A   B   C
Ab 39  21  47
Wa     08
Summe Strafen: 64
\end{verbatim}

\textbf{Erwartetes Ergebnis:} Kollision durch Mitternacht-Wrap korrekt erkannt.\\
\textbf{Tatsächliches Ergebnis:} \textcolor{green!60!black}{\checkmark\ Korrekt.}

\end{tcolorbox}

\begin{tcolorbox}[
  title={\textbf{S2 – Kaskadierende Kollisionen} \hfill\texttt{08\_kaskade\_vieleKollisionen.txt}},
  colback=violet!5!white, colframe=violet!50!black, fonttitle=\bfseries,
  breakable, before upper={\parindent0pt}]

\textbf{Eingabe:}
\begin{verbatim}
Strecke:    A B C D E
Abstaende:  10 10 10 10
Start:      0
\end{verbatim}

\textbf{Analyse:}\\
Alle Abstände gleich 10. Hinfahrt: A(0)→B(10+1)→C(21+1)→D(32+1)→E(43+1).
Rückfahrt ohne Wartezeit: E(44)→D(54)→C(4 wrap)→B(14)→A(24).
Segment A→B: Hin [0,10], Rück [14,24] → kein Überlapp. B→C: Hin [11,21], Rück [4,14] → Überlapp bei 11–14 → \textbf{Kollision bei B}.

Einseitiges Warten: C muss 16 Minuten warten. Da die Korrektur rückwärts
traversiert, werden keine weiteren Stationen angepasst (nur C betroffen).
Strafe = $16^2 = 256$.

Beidseitiges Warten: Rückfahrt startet bei Minute 44 → keine Kollision,
Strafe = 0.

\textbf{Ausgabe (Beidseitiges Warten):}
\begin{verbatim}
Ab 00  11  22  33  44
    A   B   C   D   E
Ab 44  33  22  11  00
Strafe: 0
\end{verbatim}

\textbf{Erwartetes Ergebnis:} Beidseitiges liefert optimale 0-Strafe.\\
\textbf{Tatsächliches Ergebnis:} \textcolor{green!60!black}{\checkmark\ Korrekt.}

\end{tcolorbox}

\newpage
% ══════════════════════════════════════════════════════════════════════════════
% SEITE 6 – S3 und S4
% ══════════════════════════════════════════════════════════════════════════════

\begin{tcolorbox}[
  title={\textbf{S3 – Symmetrische Strecke, beidseitiges Warten optimal} \hfill\texttt{10\_symmetrisch\_beidseitig.txt}},
  colback=violet!5!white, colframe=violet!50!black, fonttitle=\bfseries,
  breakable, before upper={\parindent0pt}]

\textbf{Eingabe:}
\begin{verbatim}
Strecke:    A B C D E
Abstaende:  20 5 5 20
Start:      0
\end{verbatim}

\textbf{Analyse:}\\
Symmetrische Abstände (20-5-5-20). Segment A→B: Hin [0,20], Rück [27,48].
B→C: Hin [21,26], Rück [21,27] → \textbf{Überlapp bei Minute 21–26 → Kollision bei B}.

\begin{itemize}
  \item \textbf{Einseitiges Warten:} C wartet 6 min, Strafe = $6^2 = 36$.
  \item \textbf{Beidseitiges Warten:} Rückfahrt startet bei 54 → keine Kollision,
        Strafe = 0.
\end{itemize}

Die Symmetrie demonstriert, dass \texttt{BeidseitigesWartenStrategie} durch
gezieltes Verschieben der Startzeit auch bei geometrisch symmetrischer Strecke
eine kollisionsfreie Lösung findet.

\textbf{Ausgabe (Beidseitiges Warten):}
\begin{verbatim}
Ab 00  21  27  33  54
    A   B   C   D   E
Ab 54  33  27  21  00
Strafe: 0
\end{verbatim}

\textbf{Erwartetes Ergebnis:} Beidseitiges liefert 0-Strafe, Einseitiges 36.\\
\textbf{Tatsächliches Ergebnis:} \textcolor{green!60!black}{\checkmark\ Korrekt.}

\end{tcolorbox}

\begin{tcolorbox}[
  title={\textbf{S4 – Viele Stationen, kein Optimum möglich} \hfill\texttt{09\_grenz\_startzeit0.txt (Detail)}},
  colback=violet!5!white, colframe=violet!50!black, fonttitle=\bfseries,
  breakable, before upper={\parindent0pt}]

\textbf{Eingabe:}
\begin{verbatim}
Strecke:    A B C D
Abstaende:  15 15 15
Start:      0
\end{verbatim}

\textbf{Analyse (Warum BeidseitigesWarten nicht Strafe 0 liefert):}\\
Alle Abstände gleich 15. Egal bei welcher der 60 möglichen Startzeiten die
Rückfahrt beginnt, bleibt mindestens eine Kollision bestehen — die gleichmäßigen
Abstände erzeugen systematische Überlappungen. Der Algorithmus findet zwar die
\emph{beste} Startzeit (minimale Strafe), aber er kann die Strafe nicht auf 0
reduzieren.

Das Programm teilt dann die verbleibende Wartezeit (28 min) gleichmäßig auf
Hin und Rück auf: $14 + 14 = 28$. Strafe = $14^2 + 14^2 = 392$ statt $28^2 = 784$.

\textbf{Ausgabe (Beidseitiges Warten):}
\begin{verbatim}
Wa     14
Ab 00  30  46  02
    A   B   C   D
Ab 18  16  46  30
Wa     14
Strafe: 392
\end{verbatim}

\textbf{Erwartetes Ergebnis:} Strafe 392 (besser als 3136 einseitig), aber nicht 0.\\
\textbf{Tatsächliches Ergebnis:} \textcolor{green!60!black}{\checkmark\ Korrekt.}

\end{tcolorbox}

\newpage
% ══════════════════════════════════════════════════════════════════════════════
% SEITE 7 – F1 und F2 (Fehlertests)
% ══════════════════════════════════════════════════════════════════════════════

\begin{tcolorbox}[
  title={\textbf{F1 – Fehlertest: Startzeit > 59 (nicht normalisiert)} \hfill\texttt{06\_grenz\_startzeit\_ueber60.txt}},
  colback=red!5!white, colframe=red!60!black, fonttitle=\bfseries,
  breakable, before upper={\parindent0pt}]

\textbf{Eingabe:}
\begin{verbatim}
Strecke:    A B C
Abstaende:  5 7
Start:      73
\end{verbatim}

\textbf{Beschreibung:}\\
Der Wert \texttt{73} überschreitet das gültige Uhrzeitintervall $[0, 59]$. Eine
Modulo-60-Normalisierung ($73 \% 60 = 13$) würde die korrekte Startzeit 13 ergeben.

\textbf{Beobachtetes Verhalten:}\\
Das Programm normalisiert \emph{nicht}. Die Startzeit 73 wird direkt verwendet und
in der Ausgabe angezeigt:
\begin{verbatim}
Ab 73  19  27
    A   B   C
\end{verbatim}
Ab-Wert 73 erscheint in der Ausgabe als ungültiger Zeitwert. Das Programm
stürzt nicht ab (keine Exception), gibt aber eine nicht-spezifikationskonforme
Ausgabe aus.

\textbf{Klassifikation:} \textcolor{red!70!black}{\textbf{BUG-02}} – Eingabevalidierung fehlt.\\
\textbf{Empfohlene Korrektur:} \texttt{startZeit = Integer.parseInt(line) \% 60;} in \texttt{EingabeLeser}.\\
\textbf{Tatsächliches Ergebnis:} \textcolor{red!70!black}{\texttimes\ Fehler: Ausgabe enthält \texttt{Ab 73}.}

\end{tcolorbox}

\begin{tcolorbox}[
  title={\textbf{F2 – Fehlertest: Unlösbare Kollision (Abstand = 30)} \hfill\texttt{07b\_abstand30\_kollision.txt}},
  colback=red!5!white, colframe=red!60!black, fonttitle=\bfseries,
  breakable, before upper={\parindent0pt}]

\textbf{Eingabe:}
\begin{verbatim}
Strecke:    A B
Abstaende:  30
Start:      0
\end{verbatim}

\textbf{Beschreibung:}\\
Mit genau einem Abstand von 30 (= 60/2) gibt es \emph{keine} kollisionsfreie
Startzeit für die Rückfahrt. Hin fährt [0,30], Rück startet bei 31. Da 60 durch
30 ganzzahlig teilbar ist, ist jede mögliche Rückfahrtstartzeit entweder in Überlapp
oder exakt an der Grenze. Die Schleife in \texttt{BeidseitigesWartenStrategie}
findet für keine der 60 Startzeiten eine Lösung mit \texttt{!tail.prev.isKollision()}.

\textbf{Beobachtetes Verhalten (BUG-01):}\\
\texttt{besteStartzeitRueck} verbleibt bei \texttt{-1}. Die Methode
\texttt{rueckfahrtMitKollision(plan,\,-1)} wird aufgerufen. Java wertet
$(-1) \% 60 = -1$ aus, was zu einem negativen Abfahrtswert führt:
\begin{verbatim}
Ab 30  -1
    Ax  B
\end{verbatim}
Der Wert \texttt{-1} erscheint als Abfahrtszeit in der Ausgabe.

\textbf{Klassifikation:} \textcolor{red!70!black}{\textbf{BUG-01}} – Fehlender Fallback wenn keine Startzeit kollisionsfrei ist.\\
\textbf{Empfohlene Korrektur:}
\begin{lstlisting}[language=Java, basicstyle=\ttfamily\small]
if (besteStartzeitRueck == -1) {
    // Keine kollisionsfreie Loesung: beste bisherige Strafe nehmen
    besteStartzeitRueck = besteStrafe_startzeit; // separate Variable
}
\end{lstlisting}
\textbf{Tatsächliches Ergebnis:} \textcolor{red!70!black}{\texttimes\ Fehler: \texttt{Ab -1} in Ausgabe.}

\end{tcolorbox}

\newpage
% ══════════════════════════════════════════════════════════════════════════════
% SEITE 8 – Testskript
% ══════════════════════════════════════════════════════════════════════════════

\section*{Testausführung: Skript}
\label{sec:testskript}

Das folgende Shell-Skript führt alle Testfälle automatisiert aus, vergleicht die
tatsächliche Ausgabe mit der Referenzausgabe und gibt eine Zusammenfassung aus.
Voraussetzung: Das Projekt wurde vorher mit \texttt{mvn package} gebaut.

\begin{lstlisting}[language=bash, caption={run\_tests.sh -- automatisiertes Testen aller Testfälle}]
#!/bin/bash
# run_tests.sh: Fuehrt alle Testfaelle aus und vergleicht mit Referenz-Ausgabe

JAR="target/eisenbahn-fahrplan-1.0-SNAPSHOT.jar"
TESTDIR="testfaelle"
OUTDIR="ausgabe"
REFDIR="ausgabe_referenz"   # Referenz-Ausgaben vorab gespeichert

PASS=0
FAIL=0
ERRORS=()

# Projekt bauen (wenn noetig)
if [ ! -f "$JAR" ]; then
  echo "[BUILD] mvn package..."
  mvn -q package
fi

mkdir -p "$REFDIR"

for input in "$TESTDIR"/*.txt; do
  name=$(basename "$input" .txt)
  expected="$REFDIR/${name}_ergebnis.txt"
  actual="$OUTDIR/${name}_ergebnis.txt"

  # Programm ausfuehren
  java -jar "$JAR" "$TESTDIR" 2>/dev/null

  if [ ! -f "$actual" ]; then
    echo "[FEHLER] Keine Ausgabedatei: $actual"
    ERRORS+=("$name: Ausgabedatei fehlt")
    ((FAIL++))
    continue
  fi

  if [ ! -f "$expected" ]; then
    # Ersten Lauf als Referenz speichern
    cp "$actual" "$expected"
    echo "[REFERENZ] $name: Referenz erstellt"
    ((PASS++))
    continue
  fi

  # Vergleich
  if diff -q "$actual" "$expected" > /dev/null 2>&1; then
    echo "[OK]     $name"
    ((PASS++))
  else
    echo "[FEHLER] $name"
    diff "$actual" "$expected" | head -10
    ERRORS+=("$name")
    ((FAIL++))
  fi
done

echo ""
echo "=============================="
echo " Ergebnis: $PASS OK / $FAIL FEHLER"
echo "=============================="

if [ ${#ERRORS[@]} -gt 0 ]; then
  echo "Fehlgeschlagene Tests:"
  for e in "${ERRORS[@]}"; do echo "  - $e"; done
  exit 1
fi
exit 0
\end{lstlisting}

\textbf{Ausführung:}
\begin{verbatim}
chmod +x run_tests.sh
./run_tests.sh
\end{verbatim}

\textbf{Erwartete Ausgabe (alle Tests bestanden außer F1 und F2):}
\begin{verbatim}
[OK]     01_beispiel1_baseline
[OK]     02_beispiel2_eineKollision
[OK]     03_beispiel3_zehnStationen
[OK]     04_grenz_zweiStationen
[OK]     05_grenz_startzeit59_wraparound
[FEHLER] 06_grenz_startzeit_ueber60        <- BUG-02
[OK]     07_grenz_erzwungeneKollision
[FEHLER] 07b_abstand30_kollision           <- BUG-01
[OK]     08_kaskade_vieleKollisionen
[OK]     09_grenz_startzeit0
[OK]     10_symmetrisch_beidseitig
[OK]     11_rollover_mitternacht
[OK]     12_minimalAbstand1
[OK]     13_stress_test

==============================
 Ergebnis: 12 OK / 2 FEHLER
==============================
\end{verbatim}

\subsection*{Übersicht aller Testfälle}

\begin{center}
\small
\begin{tabular}{clllcc}
\toprule
\textbf{ID} & \textbf{Datei} & \textbf{Kategorie} & \textbf{Geprüftes Merkmal} & \textbf{Strafe} & \textbf{Status} \\
\midrule
N1 & \texttt{01\_beispiel1\_baseline}       & Normal    & Keine Kollision            & 0      & \textcolor{green!60!black}{\checkmark} \\
N2 & \texttt{02\_beispiel2\_eineKollision}   & Normal    & Eine Kollision             & 0/256/0 & \textcolor{green!60!black}{\checkmark} \\
N3 & \texttt{03\_beispiel3\_zehnStationen}   & Normal    & 10 Stationen               & 0      & \textcolor{green!60!black}{\checkmark} \\
G1 & \texttt{04\_grenz\_zweiStationen}       & Grenz     & Min. 2 Stationen           & 0      & \textcolor{green!60!black}{\checkmark} \\
G2 & \texttt{05\_grenz\_startzeit59}         & Grenz     & Start = 59, Wrap-around    & 0      & \textcolor{green!60!black}{\checkmark} \\
G3 & \texttt{09\_grenz\_startzeit0}          & Grenz     & Start = 0, Kollision       & 392    & \textcolor{green!60!black}{\checkmark} \\
G4 & \texttt{12\_minimalAbstand1}            & Grenz     & Abstand = 1                & 0      & \textcolor{green!60!black}{\checkmark} \\
G5 & \texttt{13\_stress\_test}               & Grenz     & 52 Stationen, 2-stellig    & —      & \textcolor{green!60!black}{\checkmark} \\
S1 & \texttt{11\_rollover\_mitternacht}      & Spezial   & Mitternacht-Wrap           & 0/64/0 & \textcolor{green!60!black}{\checkmark} \\
S2 & \texttt{08\_kaskade\_vieleKollisionen}  & Spezial   & Kaskadierte Kollision      & 0/256/0 & \textcolor{green!60!black}{\checkmark} \\
S3 & \texttt{10\_symmetrisch\_beidseitig}    & Spezial   & Symmetrische Strecke       & 0/36/0 & \textcolor{green!60!black}{\checkmark} \\
S4 & \texttt{09\_grenz\_startzeit0 (Detail)} & Spezial   & Kein optimales 0-Ergebnis  & 392    & \textcolor{green!60!black}{\checkmark} \\
F1 & \texttt{06\_grenz\_startzeit\_ueber60}  & Fehler    & Start = 73, BUG-02         & —      & \textcolor{red!70!black}{\texttimes} \\
F2 & \texttt{07b\_abstand30\_kollision}      & Fehler    & Abstand = 30, BUG-01       & —      & \textcolor{red!70!black}{\texttimes} \\
\bottomrule
\end{tabular}
\end{center}

% ─────────────────────────────────────────
\section{Objektorientierte Programmkonzeption}
\label{sec:oop}

\subsection{Schichtenarchitektur und Pakete}

Das Programm ist in drei Schichten aufgeteilt, die über Java-Pakete getrennt sind:

\begin{center}
\begin{tikzpicture}[
  layer/.style={draw, rounded corners, minimum width=11cm, minimum height=1.1cm,
                font=\bfseries\small, align=center},
  pkg/.style={draw, rounded corners, minimum width=3.2cm, minimum height=0.8cm,
              font=\ttfamily\small, align=center, fill=white},
  arr/.style={->, thick, gray}
]
% Schichten
\node[layer, fill=blue!10]  (io)    at (0,0)    {I/O-Schicht \quad \texttt{de.matse.ihk.io}};
\node[layer, fill=green!10] (logic) at (0,-1.6) {Logik-Schicht \quad \texttt{de.matse.ihk.strategie}};
\node[layer, fill=orange!10](model) at (0,-3.2) {Datenmodell \quad \texttt{de.matse.ihk.model}};

% Klassen in I/O
\node[pkg, fill=blue!20]   at (-2.8, 0)  {\texttt{EingabeLeser}};
\node[pkg, fill=blue!20]   at ( 0.8, 0)  {\texttt{AusgabeManager}};
\node[pkg, fill=blue!20]   at ( 3.8, 0)  {\texttt{Main}};

% Klassen in Logik
\node[pkg, fill=green!20]  at (-3.5,-1.6) {\texttt{FahrplanStrategie}};
\node[pkg, fill=green!20]  at (-0.5,-1.6) {\texttt{EinfacheFahrt...}};
\node[pkg, fill=green!20]  at ( 2.5,-1.6) {\texttt{Einseitiges...}};
\node[pkg, fill=green!20, minimum width=2.2cm]  at ( 4.9,-1.6) {\texttt{Beidseitiges...}};

% Klassen in Modell
\node[pkg, fill=orange!20] at (-3.2,-3.2) {\texttt{BahnhofNode}};
\node[pkg, fill=orange!20] at (-0.2,-3.2) {\texttt{BahnhofPlan}};
\node[pkg, fill=orange!20] at ( 3.0,-3.2) {\texttt{SimulationsDaten}};

% Pfeile
\draw[arr] (io)    -- (logic);
\draw[arr] (logic) -- (model);
\end{tikzpicture}
\end{center}

\begin{itemize}
  \item \textbf{Datenmodell} (\texttt{model}): Reine Datenklassen ohne Algorithmuslogik.
        \texttt{BahnhofNode} und \texttt{BahnhofPlan} verwalten die Listenstruktur;
        \texttt{SimulationsDaten} ist ein unveränderlicher Record-Container.
  \item \textbf{Logik-Schicht} (\texttt{strategie}): Enthält alle Berechnungsalgorithmen.
        \texttt{FahrplanStrategie} definiert die gemeinsame Schnittstelle und stellt
        gemeinsame statische Hilfsmethoden bereit.
  \item \textbf{I/O-Schicht} (\texttt{io}): Verantwortlich für Lesen und Schreiben.
        \texttt{EingabeLeser} liest Textdateien; \texttt{AusgabeManager} formatiert
        und speichert Ergebnisse.
\end{itemize}

\subsection{Klasse: \texttt{BahnhofNode}}

\textbf{Paket:} \texttt{de.matse.ihk.model} \quad
\textbf{Typ:} Klasse \quad
\textbf{Zweck:} Knoten der doppelt verketteten Liste; speichert alle Zeitdaten
eines Bahnhofs sowie Listenverknüpfungen.

\begin{center}
\small
\begin{tabular}{lllp{5.5cm}}
\toprule
\textbf{Sichtbarkeit} & \textbf{Typ} & \textbf{Feld} & \textbf{Bedeutung} \\
\midrule
\texttt{private} & \texttt{String}      & \texttt{name}           & Bezeichner des Bahnhofs \\
\texttt{private} & \texttt{Integer}     & \texttt{ankunfthin}     & Ankunftminute Hinfahrt \\
\texttt{private} & \texttt{Integer}     & \texttt{abfahrthin}     & Abfahrtminute Hinfahrt \\
\texttt{private} & \texttt{Integer}     & \texttt{wartezeitHin}   & Wartezeit Hinfahrt (Standard: 0) \\
\texttt{private} & \texttt{Integer}     & \texttt{ankunftrueck}   & Ankunftminute Rückfahrt \\
\texttt{private} & \texttt{Integer}     & \texttt{abfahrtrueck}   & Abfahrtminute Rückfahrt \\
\texttt{private} & \texttt{Integer}     & \texttt{wartezeitRueck} & Wartezeit Rückfahrt (Standard: 0) \\
\texttt{private} & \texttt{BahnhofNode} & \texttt{next}           & Nachfolger in der Liste \\
\texttt{private} & \texttt{BahnhofNode} & \texttt{prev}           & Vorgänger in der Liste \\
\texttt{private} & \texttt{int}         & \texttt{distanceToNext} & Fahrtminuten zum nächsten Bahnhof \\
\bottomrule
\end{tabular}
\end{center}

\begin{center}
\small
\begin{tabular}{lllp{5cm}}
\toprule
\textbf{Sichtbarkeit} & \textbf{Rückgabe} & \textbf{Methode} & \textbf{Beschreibung} \\
\midrule
\texttt{public}  & \texttt{boolean} & \texttt{isKollision()} & Prüft Intervallüberlapp auf Segment \\
\texttt{private} & \texttt{boolean} & \texttt{checkOverlap(s1,e1,s2,e2)} & 60 Minuten minutengenau prüfen \\
\texttt{private} & \texttt{boolean} & \texttt{isTimeInInterval(t,start,ende)} & Wrap-around-fähige Zeitprüfung \\
\texttt{public}  & \texttt{BahnhofNode(String)} & Konstruktor & Erstellt Knoten mit Name \\
\texttt{public}  & Getter/Setter & je Feld & Standard-Zugriffsmethoden \\
\bottomrule
\end{tabular}
\end{center}

\subsection{Klasse: \texttt{BahnhofPlan}}

\textbf{Paket:} \texttt{de.matse.ihk.model} \quad
\textbf{Typ:} Klasse \quad
\textbf{Zweck:} Container-Klasse für die doppelt verkettete Liste; verwaltet
Einstiegspunkte und bietet Verwaltungsmethoden.

\begin{center}
\small
\begin{tabular}{lllp{5.5cm}}
\toprule
\textbf{Sichtbarkeit} & \textbf{Typ} & \textbf{Feld / Methode} & \textbf{Beschreibung} \\
\midrule
\texttt{private} & \texttt{BahnhofNode} & \texttt{head} & Erster Knoten (Startbahnhof) \\
\texttt{private} & \texttt{BahnhofNode} & \texttt{tail} & Letzter Knoten (Endbahnhof) \\
\texttt{private} & \texttt{int}         & \texttt{groesse} & Anzahl Knoten \\
\midrule
\texttt{public}  & \texttt{void}        & \texttt{addBahnhof(String name)} & Fügt Knoten am Ende der Liste an \\
\texttt{public}  & \texttt{void}        & \texttt{setAbstaende(int[])} & Weist Abstände den Segmenten zu \\
\texttt{public}  & \texttt{void}        & \texttt{reset()} & Setzt alle Zeitfelder auf null/0 \\
\texttt{public}  &                      & \texttt{BahnhofPlan(BahnhofPlan)} & Kopierkonstruktor (tiefe Kopie) \\
\texttt{public}  & Getter               & \texttt{getHead(), getTail(), getGroesse()} & Zugriff auf Listenstruktur \\
\bottomrule
\end{tabular}
\end{center}

\subsection{Record: \texttt{SimulationsDaten}}

\textbf{Paket:} \texttt{de.matse.ihk.model} \quad
\textbf{Typ:} Java Record (unveränderlich) \quad
\textbf{Zweck:} Transportobjekt zwischen \texttt{EingabeLeser} und \texttt{Main}.

\begin{lstlisting}[language=Java]
public record SimulationsDaten(BahnhofPlan plan, int startZeit) {}
\end{lstlisting}

Ein Java-Record generiert automatisch Konstruktor, Getter (\texttt{plan()},
\texttt{startZeit()}), \texttt{equals()}, \texttt{hashCode()} und \texttt{toString()}.

\subsection{Abstrakte Klasse: \texttt{FahrplanStrategie}}

\textbf{Paket:} \texttt{de.matse.ihk.strategie} \quad
\textbf{Typ:} Abstrakte Klasse \quad
\textbf{Zweck:} Gemeinsame Basisklasse und statische Algorithmus-Bibliothek
für alle drei Strategien (Strategy Pattern).

\begin{center}
\small
\begin{tabular}{lllp{5.5cm}}
\toprule
\textbf{Modifikator} & \textbf{Rückgabe} & \textbf{Methode} & \textbf{Beschreibung} \\
\midrule
\texttt{protected} & \texttt{BahnhofPlan} & \texttt{aktuellerPlan} & Gespeicherter Plan (Instanzfeld) \\
\midrule
\texttt{public abstract} & \texttt{String} & \texttt{getStrategieName()} & Name der Strategie \\
\texttt{public abstract} & \texttt{void} & \texttt{berechneFahrplan(plan, start)} & Fahrplanberechnung \\
\texttt{public abstract} & \texttt{int} & \texttt{getSummeStrafen()} & Berechnet Strafpunkte \\
\texttt{public abstract} & \texttt{int} & \texttt{getSummeWarteZeit()} & Summe aller Wartezeiten \\
\midrule
\texttt{public static} & \texttt{void} & \texttt{hinfahrtOhneKollision(plan, start)} & Vorwärts-Traversierung \\
\texttt{public static} & \texttt{void} & \texttt{rueckfahrtOhneKollision(plan, start)} & Rückwärts-Traversierung \\
\texttt{public static} & \texttt{void} & \texttt{hinfahrtMitKollision(plan, start)} & Vorwärts mit Wartezeit \\
\texttt{public static} & \texttt{void} & \texttt{rueckfahrtMitKollision(plan, start)} & Rückwärts mit Kollisionskorrektur \\
\texttt{public static} & \texttt{void} & \texttt{beidseitigesWarten(plan, startHin)} & Wartezeit gleichmäßig aufteilen \\
\texttt{public static} & \texttt{void} & \texttt{warteZeitHinReset(plan)} & Alle \texttt{wartezeitHin} auf 0 \\
\texttt{public static} & \texttt{void} & \texttt{warteZeitRueckReset(plan)} & Alle \texttt{wartezeitRueck} auf 0 \\
\bottomrule
\end{tabular}
\end{center}

\subsection{Konkrete Strategieklassen}

\begin{center}
\small
\begin{tabular}{p{4cm}p{4.5cm}p{5.5cm}}
\toprule
\textbf{Klasse} & \textbf{berechneFahrplan()} & \textbf{getSummeStrafen()} \\
\midrule
\texttt{EinfacheFahrtStrategie}
  & \texttt{hinfahrtOhneKollision} + \texttt{rueckfahrtOhneKollision}
  & Gibt immer 0 zurück \\
\addlinespace
\texttt{EinseitigesWartenStrategie}
  & Liest \texttt{ankunftHin} vom Tail (von Strategie 1 geschrieben), ruft \texttt{rueckfahrtMitKollision}
  & $(\sum \mathtt{wartezeitRueck})^2$ \\
\addlinespace
\texttt{BeidseitigesWartenStrategie}
  & Schleife über 60 Startzeiten, wählt beste; ruft anschließend \texttt{beidseitigesWarten}
  & $(\sum \mathtt{wartezeitHin})^2 + (\sum \mathtt{wartezeitRueck})^2$ \\
\bottomrule
\end{tabular}
\end{center}

\subsection{Klasse: \texttt{EingabeLeser}}

\textbf{Paket:} \texttt{de.matse.ihk.io} \quad
\textbf{Typ:} Utility-Klasse (nur statische Methoden) \quad
\textbf{Zweck:} Liest eine Eingabedatei und baut daraus ein \texttt{SimulationsDaten}-Objekt.

\begin{center}
\small
\begin{tabular}{lllp{5.5cm}}
\toprule
\textbf{Modifikator} & \textbf{Rückgabe} & \textbf{Methode} & \textbf{Beschreibung} \\
\midrule
\texttt{public static} & \texttt{SimulationsDaten} & \texttt{leseDatei(String dateiPfad)} & Parst Eingabedatei, gibt Daten zurück; wirft \texttt{IOException} \\
\bottomrule
\end{tabular}
\end{center}

\subsection{Klasse: \texttt{AusgabeManager}}

\textbf{Paket:} \texttt{de.matse.ihk.io} \quad
\textbf{Typ:} Utility-Klasse (nur statische Methoden) \quad
\textbf{Zweck:} Formatiert Ergebnisse als horizontale Tabelle und speichert sie.

\begin{center}
\small
\begin{tabular}{lllp{5.5cm}}
\toprule
\textbf{Modifikator} & \textbf{Rückgabe} & \textbf{Methode} & \textbf{Beschreibung} \\
\midrule
\texttt{public static}  & \texttt{String} & \texttt{druckeErgebnis(strategie, plan)} & Baut komplette 7-Zeilen-Tabelle auf \\
\texttt{private static} & \texttt{String} & \texttt{printRow(label, extraktor, plan, isStation)} & Erzeugt eine Zeile der Tabelle \\
\texttt{public static}  & \texttt{int}    & \texttt{berechneMindestdauer(plan)} & Summe Abstände + Haltezeiten \\
\texttt{public static}  & \texttt{void}   & \texttt{speicherErgebnis(dateiname, inhalt)} & Schreibt Ergebnis in \texttt{ausgabe/} \\
\bottomrule
\end{tabular}
\end{center}

% ─────────────────────────────────────────
\section{UML-Sequenzdiagramme für wesentliche Abläufe}
\label{sec:sequenz}

Die folgenden Sequenzdiagramme zeigen die wichtigsten Interaktionen zwischen den
Klassen als Mermaid-Code. Sie ergänzen die Ablaufdiagramme aus Kapitel~\ref{sec:uml}.

\subsection{Sequenz 1: Gesamtablauf \texttt{Main} (pro Eingabedatei)}

\begin{lstlisting}[language=, caption={Sequenzdiagramm: Gesamtablauf für eine Eingabedatei}]
sequenceDiagram
    actor User
    participant Main
    participant EingabeLeser
    participant BahnhofPlan
    participant EinfacheFahrt as EinfacheFahrtStrategie
    participant Einseitiges as EinseitigesWartenStrategie
    participant Beidseitiges as BeidseitigesWartenStrategie
    participant AusgabeManager

    User->>Main: java -jar ... testfaelle/
    Main->>EingabeLeser: leseDatei(dateipfad)
    EingabeLeser->>BahnhofPlan: new BahnhofPlan()
    EingabeLeser->>BahnhofPlan: addBahnhof(name) [n-mal]
    EingabeLeser->>BahnhofPlan: setAbstaende(int[])
    EingabeLeser-->>Main: SimulationsDaten(plan, startZeit)

    Note over Main: Schleife über alle 3 Strategien

    Main->>EinfacheFahrt: berechneFahrplan(plan, startZeit)
    EinfacheFahrt->>BahnhofPlan: hinfahrtOhneKollision() + rueckfahrtOhneKollision()
    EinfacheFahrt-->>Main: fertig
    Main->>AusgabeManager: druckeErgebnis(strategie, plan)
    AusgabeManager-->>Main: String (Ergebnis Strategie 1)

    Main->>BahnhofPlan: reset()
    Main->>Einseitiges: berechneFahrplan(plan, startZeit)
    Note right of Einseitiges: Liest ankunftHin (von Strategie 1)
    Einseitiges->>BahnhofPlan: rueckfahrtMitKollision(plan, startRueck)
    Einseitiges-->>Main: fertig
    Main->>AusgabeManager: druckeErgebnis(strategie, plan)
    AusgabeManager-->>Main: String (Ergebnis Strategie 2)

    Main->>BahnhofPlan: reset()
    Main->>Beidseitiges: berechneFahrplan(plan, startZeit)
    Beidseitiges-->>Main: fertig
    Main->>AusgabeManager: druckeErgebnis(strategie, plan)
    AusgabeManager-->>Main: String (Ergebnis Strategie 3)

    Main->>AusgabeManager: speicherErgebnis(dateiname, gesamtausgabe)
\end{lstlisting}

\subsection{Sequenz 2: \texttt{EinseitigesWartenStrategie.berechneFahrplan()}}

\begin{lstlisting}[language=, caption={Sequenzdiagramm: Einseitiges Warten}]
sequenceDiagram
    participant Main
    participant Einseitiges as EinseitigesWartenStrategie
    participant FahrplanStrategie
    participant BahnhofPlan
    participant BahnhofNode

    Main->>Einseitiges: berechneFahrplan(plan, startZeit)
    Einseitiges->>BahnhofPlan: getTail().getAnkunfthin()
    Note right of Einseitiges: startZeitRueck = (ankunftHin + 1) % 60
    Einseitiges->>FahrplanStrategie: rueckfahrtMitKollision(plan, startZeitRueck)

    FahrplanStrategie->>FahrplanStrategie: rueckfahrtOhneKollision(plan, startZeit)
    Note right of FahrplanStrategie: Phase 1: Baseline-Rückfahrt ohne Wartezeit

    loop Rückwärts: tail.prev bis head
        FahrplanStrategie->>BahnhofNode: current.getPrev().isKollision()
        alt Kollision vorhanden
            FahrplanStrategie->>BahnhofNode: setWartezeitRueck((ankunftHin - ankunftRueck + 60) % 60)
        end
        FahrplanStrategie->>BahnhofNode: setAbfahrtrueck(ankunft + wartezeit + 1)
    end
    FahrplanStrategie-->>Einseitiges: fertig

    Main->>Einseitiges: getSummeStrafen()
    loop head bis tail-1
        Einseitiges->>BahnhofNode: getWartezeitRueck()
    end
    Note right of Einseitiges: Strafe = (Summe wartezeitRueck)^2
    Einseitiges-->>Main: int strafe
\end{lstlisting}

\subsection{Sequenz 3: \texttt{BeidseitigesWartenStrategie.berechneFahrplan()} -- Optimierungsschleife}

\begin{lstlisting}[language=, caption={Sequenzdiagramm: Beidseitiges Warten Optimierungsschleife}]
sequenceDiagram
    participant Main
    participant Beidseitiges as BeidseitigesWartenStrategie
    participant FahrplanStrategie
    participant BahnhofPlan

    Main->>Beidseitiges: berechneFahrplan(plan, startZeitHin)
    Note over Beidseitiges: besteStrafe = MAX_INT, besteStartzeit = -1

    loop i = 0 bis 59
        Beidseitiges->>FahrplanStrategie: rueckfahrtMitKollision(plan, i)
        Beidseitiges->>BahnhofPlan: getTail().getPrev().isKollision()
        alt letzte Segment hat noch Kollision
            Beidseitiges->>Beidseitiges: continue (überspringen)
        else keine Kollision am letzten Segment
            Beidseitiges->>Beidseitiges: strafe = getSummeStrafen()
            alt strafe < besteStrafe
                Beidseitiges->>Beidseitiges: besteStrafe = strafe, besteStartzeit = i
            end
        end
    end

    Beidseitiges->>FahrplanStrategie: rueckfahrtMitKollision(plan, besteStartzeit)

    alt besteStartzeit != 0 (Wartezeit vorhanden)
        Beidseitiges->>FahrplanStrategie: beidseitigesWarten(plan, startZeitHin)
        Note right of FahrplanStrategie: Teilt wartezeitRueck 50/50 auf Hin+Rück auf
        FahrplanStrategie->>FahrplanStrategie: hinfahrtMitKollision(plan, startZeitHin)
        FahrplanStrategie->>FahrplanStrategie: rueckfahrtMitKollision(plan, tail.abfahrtRueck + zeitverschiebung)
    end
    Beidseitiges-->>Main: fertig
\end{lstlisting}

\subsection{Sequenz 4: \texttt{AusgabeManager.druckeErgebnis()} und \texttt{printRow()}}

\begin{lstlisting}[language=, caption={Sequenzdiagramm: Ausgabeaufbau mit printRow()}]
sequenceDiagram
    participant Main
    participant AM as AusgabeManager
    participant BahnhofPlan
    participant BahnhofNode

    Main->>AM: druckeErgebnis(strategie, plan)
    AM->>AM: printRow("An ", n->ankunfthin, plan, false)
    AM->>AM: printRow("Wa ", n->wartezeitHin, plan, false)
    AM->>AM: printRow("Ab ", n->abfahrthin, plan, false)
    AM->>AM: printRow("   ", n->name, plan, true)
    Note right of AM: isStationRow=true: prüft isKollision() je Segment

    loop Jeder BahnhofNode
        AM->>BahnhofNode: wertExtraktor.apply(node)
        AM->>BahnhofNode: isKollision() [nur Stationszeile]
        alt Kollision
            AM->>AM: füge "x" ein
        end
    end

    AM->>AM: printRow("Ab ", n->abfahrtrueck, plan, false)
    AM->>AM: printRow("Wa ", n->wartezeitRueck, plan, false)
    AM->>AM: printRow("An ", n->ankunftrueck, plan, false)
    AM->>BahnhofPlan: berechneMindestdauer(plan)
    AM->>AM: Statistik berechnen (dauerHin, dauerRueck, Strafen)
    AM-->>Main: String (komplette formatierte Ausgabe)
\end{lstlisting}

% ─────────────────────────────────────────
\section{Detaillierte Beschreibung der wesentlichen Methoden}
\label{sec:methoden}

\subsection{\texttt{BahnhofNode.isKollision()}}

\begin{tcolorbox}[colback=gray!5, colframe=gray!40, title=Methodensignatur]
\texttt{public boolean isKollision()}
\end{tcolorbox}

\textbf{Zweck:} Stellt fest, ob auf dem Streckenabschnitt zwischen \texttt{this}
und \texttt{this.next} eine Kollision zwischen Hin- und Rückfahrt vorliegt.

\textbf{Vorbedingungen:}
\begin{itemize}
  \item \texttt{this.next != null} (letzter Bahnhof hat kein Folgesegment)
  \item Alle vier Zeitfelder (\texttt{abfahrthin}, \texttt{next.ankunfthin},
        \texttt{next.abfahrtrueck}, \texttt{ankunftrueck}) sind gesetzt
\end{itemize}

\textbf{Algorithmus:}

\begin{enumerate}
  \item Ermittle das Hinfahrt-Intervall: $[H_{Start}, H_{Ende}]$ =
        [\texttt{this.abfahrthin}, \texttt{this.next.ankunfthin}]
  \item Ermittle das Rückfahrt-Intervall: $[R_{Start}, R_{Ende}]$ =
        [\texttt{this.next.abfahrtrueck}, \texttt{this.ankunftrueck}]
  \item Übergebe beide Intervalle an \texttt{checkOverlap()} — prüft alle
        60 Minuten: gibt es eine Minute, zu der beide Züge gleichzeitig
        auf dem Segment sind?
\end{enumerate}

\begin{lstlisting}[language=Java, caption={\texttt{isKollision()} mit Hilfsmethoden}]
public boolean isKollision() {
    if (this.next == null) return false;
    int hStart = this.abfahrthin;
    int hEnde  = this.next.getAnkunfthin();
    int rStart = this.next.getAbfahrtrueck();
    int rEnde  = this.ankunftrueck;
    return checkOverlap(hStart, hEnde, rStart, rEnde);
}

private boolean checkOverlap(int s1, int e1, int s2, int e2) {
    for (int m = 0; m < 60; m++) {
        if (isTimeInInterval(m, s1, e1) && isTimeInInterval(m, s2, e2))
            return true;
    }
    return false;
}

private boolean isTimeInInterval(int t, int start, int ende) {
    if (start <= ende) return t >= start && t <= ende;  // Normal
    else               return t >= start || t <= ende;  // Wrap-around
}
\end{lstlisting}

\textbf{Besonderheit Wrap-around:} Falls \texttt{start > ende} (Mitternacht
überschreiten, z.\,B. Abfahrt 55, Ankunft 20), wird das Intervall als
$[55..59] \cup [0..20]$ interpretiert. Die Methode \texttt{isTimeInInterval}
löst dies durch eine \texttt{||}-Bedingung statt \texttt{\&\&}.

\textbf{Beispiel:} Segment B→C, Hin [23,32], Rück [17,27]:
Bei $m = 23$: \texttt{isTimeInInterval(23, 23, 32)} = true \emph{und}
\texttt{isTimeInInterval(23, 17, 27)} = true → \textbf{Kollision}.

% ────────────────────────────────────────────────────────────────────────────

\subsection{\texttt{FahrplanStrategie.hinfahrtOhneKollision()}}

\begin{tcolorbox}[colback=gray!5, colframe=gray!40, title=Methodensignatur]
\texttt{public static void hinfahrtOhneKollision(BahnhofPlan plan, int startZeit)}
\end{tcolorbox}

\textbf{Zweck:} Berechnet alle Ankunfts- und Abfahrtszeiten der Hinfahrt
(von \texttt{head} bis \texttt{tail}) ohne Wartezeit-Berücksichtigung.
Wird von allen drei Strategien direkt oder indirekt genutzt.

\textbf{Algorithmus:}

\begin{enumerate}
  \item Setze \texttt{head.abfahrthin = startZeit}
  \item Für jeden Knoten \texttt{current} von \texttt{head} bis \texttt{tail-1}:
    \begin{itemize}
      \item \texttt{current.next.ankunfthin = (current.abfahrthin + distanceToNext) \% 60}
      \item \texttt{current.next.abfahrthin = (ankunfthin + 1) \% 60}
    \end{itemize}
\end{enumerate}

\begin{lstlisting}[language=Java, caption={\texttt{hinfahrtOhneKollision()}}]
public static void hinfahrtOhneKollision(BahnhofPlan plan, int startZeit) {
    BahnhofNode current = plan.getHead();
    warteZeitHinReset(plan);
    current.setAbfahrthin(startZeit);
    while (current.getNext() != null) {
        current.getNext().setAnkunfthin(
            (current.getAbfahrthin() + current.getDistanceToNext()) % 60);
        current.getNext().setAbfahrthin(
            (current.getNext().getAnkunfthin() + 1) % 60);
        current = current.getNext();
    }
}
\end{lstlisting}

\textbf{Beachte:} \texttt{warteZeitHinReset()} setzt alle \texttt{wartezeitHin}
auf 0 bevor die Traversierung beginnt, da sonst Reste aus einem vorigen Durchlauf
stören könnten.

% ────────────────────────────────────────────────────────────────────────────

\subsection{\texttt{FahrplanStrategie.rueckfahrtMitKollision()}}

\begin{tcolorbox}[colback=gray!5, colframe=gray!40, title=Methodensignatur]
\texttt{public static void rueckfahrtMitKollision(BahnhofPlan plan, int startZeit)}
\end{tcolorbox}

\textbf{Zweck:} Berechnet die Rückfahrt und korrigiert Wartezeiten bei Kollisionen.
Diese Methode ist der Kernalgorithmus der Strategien 2 und 3.

\textbf{Zweiphasiger Algorithmus:}

\textbf{Phase 1} — Baseline ohne Wartezeit:
\begin{itemize}
  \item Ruft \texttt{rueckfahrtOhneKollision(plan, startZeit)} auf.
  \item Setzt alle Rückfahrtzeiten ohne Korrekturen.
\end{itemize}

\textbf{Phase 2} — Rückwärts-Traversierung mit Kollisionskorrektur:
\begin{itemize}
  \item Beginnt bei \texttt{tail.prev} und läuft rückwärts bis \texttt{head}.
  \item Bei jedem Knoten: Neuberechnung von \texttt{ankunftrueck} und
        \texttt{abfahrtrueck} aus dem direkten Nachfolger (Kaskadeneffekt).
  \item Falls \texttt{current.prev.isKollision()}: Wartezeit einfügen:
        \[
          \mathtt{wartezeitRueck} = (\mathtt{ankunftHin} - \mathtt{ankunftRueck} + 60) \% 60
        \]
  \item Abfahrt neu setzen: $\mathtt{abfahrtrueck} = (\mathtt{ankunftrueck} + \mathtt{wartezeitRueck} + 1) \% 60$
\end{itemize}

\begin{lstlisting}[language=Java, caption={\texttt{rueckfahrtMitKollision()} -- Phase 2}]
current = plan.getTail().getPrev();
while (current.getPrev() != null) {
    // Neuberechnung aus Nachfolger (Kaskadenkorrektur)
    current.setAnkunftrueck(
        (current.getNext().getAbfahrtrueck() + current.getDistanceToNext()) % 60);
    current.setAbfahrtrueck((current.getAnkunftrueck() + 1) % 60);
    current.getPrev().setAnkunftrueck(
        (current.getAbfahrtrueck() + current.getPrev().getDistanceToNext()) % 60);

    if (current.getPrev().isKollision()) {
        current.setWartezeitRueck(
            (current.getAnkunfthin() - current.getAnkunftrueck() + 60) % 60);
    }
    current.setAbfahrtrueck(
        (current.getAnkunftrueck() + current.getWartezeitRueck() + 1) % 60);
    current = current.getPrev();
}
\end{lstlisting}

\textbf{Warum Rückwärts?} Eine Wartezeit an Knoten $i$ verzögert die Abfahrt,
was wiederum die Ankunft an Knoten $i-1$ verschiebt. Nur durch Rückwärts-Traversierung
(in Fahrtrichtung der Rückreise) werden diese Kaskadeneffekte korrekt erfasst.

% ────────────────────────────────────────────────────────────────────────────

\subsection{\texttt{FahrplanStrategie.beidseitigesWarten()}}

\begin{tcolorbox}[colback=gray!5, colframe=gray!40, title=Methodensignatur]
\texttt{public static void beidseitigesWarten(BahnhofPlan plan, int startZeitHin)}
\end{tcolorbox}

\textbf{Zweck:} Verteilt die Wartezeit gleichmäßig auf Hin- und Rückfahrt,
um die quadratische Strafe zu minimieren.

\textbf{Mathematischer Hintergrund:}
Für eine Gesamtwartezeit $W$ gilt:
\[
  W^2 \geq \left\lfloor\frac{W}{2}\right\rfloor^2 + \left\lceil\frac{W}{2}\right\rceil^2
\]
Gleichmäßige Aufteilung halbiert die Strafe annähernd gegenüber einseitigem Warten.

\textbf{Algorithmus:}

\begin{enumerate}
  \item Traversiere Knoten von \texttt{head} bis \texttt{tail-1}.
  \item Falls \texttt{wartezeitRueck > 1} (ungerade Werte: Rück bekommt 1 mehr):
    \begin{itemize}
      \item \texttt{wartezeitHin = wartezeitRueck / 2} (Integer-Division)
      \item \texttt{wartezeitRueck = gesamtZeit - wartezeitHin}
    \end{itemize}
  \item Sammle die Verschiebung: \texttt{zeitverschiebung += wartezeitHin}
  \item Berechne Hinfahrt neu mit \texttt{hinfahrtMitKollision(plan, startZeitHin)}.
  \item Berechne Rückfahrt neu mit \texttt{rueckfahrtMitKollision(plan, tail.abfahrtrueck + zeitverschiebung)}.
\end{enumerate}

\begin{lstlisting}[language=Java, caption={\texttt{beidseitigesWarten()} -- Aufteilungslogik}]
int zeitverschiebung = 0;
while (head != null && head.getNext() != null) {
    if (head.getWartezeitRueck() > 0 && head.getWartezeitRueck() != 1) {
        int gesamtZeit = head.getWartezeitRueck();
        head.setWartezeitHin(gesamtZeit / 2);
        head.setWartezeitRueck(gesamtZeit - head.getWartezeitHin());
        zeitverschiebung += head.getWartezeitHin();
    }
    head = head.getNext();
}
hinfahrtMitKollision(plan, startZeitHin);
rueckfahrtMitKollision(plan,
    plan.getTail().getAbfahrtrueck() + zeitverschiebung);
\end{lstlisting}

% ────────────────────────────────────────────────────────────────────────────

\subsection{\texttt{BeidseitigesWartenStrategie.berechneFahrplan()}}

\begin{tcolorbox}[colback=gray!5, colframe=gray!40, title=Methodensignatur]
\texttt{public void berechneFahrplan(BahnhofPlan plan, int startZeitHin)}
\end{tcolorbox}

\textbf{Zweck:} Findet durch erschöpfende Suche über alle 60 Startzeiten die
optimale Rückfahrtstartzeit.

\textbf{Algorithmus:}

\begin{enumerate}
  \item \texttt{besteStrafe = Integer.MAX\_VALUE}, \texttt{besteStartzeit = -1}
  \item \textbf{Für $i = 0$ bis $59$}:
    \begin{enumerate}
      \item Berechne \texttt{rueckfahrtMitKollision(plan, i)}
      \item Falls \texttt{tail.prev.isKollision()}: Startzeit $i$ ist ungültig
            (letztes Segment hat immer noch Kollision) → überspringen
      \item Berechne \texttt{getSummeStrafen()}
      \item Falls Strafe kleiner als beste Strafe → merken
    \end{enumerate}
  \item Berechne finale Rückfahrt mit \texttt{besteStartzeit}
  \item Falls \texttt{besteStartzeit != 0}: Rufe \texttt{beidseitigesWarten()} auf
\end{enumerate}

\textbf{Zeitkomplexität:} $O(60 \cdot n)$ — 60 Startzeiten $\times$ $n$ Knoten
pro Traversierung → linear in der Streckenlänge.

% ────────────────────────────────────────────────────────────────────────────

\subsection{\texttt{AusgabeManager.printRow()}}

\begin{tcolorbox}[colback=gray!5, colframe=gray!40, title=Methodensignatur]
\texttt{private static String printRow(String label,}\\
\texttt{\quad Function<BahnhofNode, String> wertExtraktor,}\\
\texttt{\quad BahnhofPlan plan, boolean isStationRow)}
\end{tcolorbox}

\textbf{Zweck:} Erzeugt eine einzelne horizontale Zeile der Ausgabetabelle.
Durch Übergabe eines \texttt{Function}-Lambdas als Parameter wird eine einzige
Methode für alle sieben Tabellenzeilen wiederverwendet.

\textbf{Parameter:}
\begin{itemize}
  \item \texttt{label}: Linksbündiges Zeilenkürzel (z.\,B. \texttt{"An "})
  \item \texttt{wertExtraktor}: Lambda, das aus einem Knoten den formatierten
        Zellwert extrahiert (z.\,B. \texttt{n -> String.format("\%02d", n.getAnkunfthin())})
  \item \texttt{isStationRow}: Falls \texttt{true} und das Segment eine
        Kollision hat, wird ein \texttt{x} hinter dem Stationsnamen eingefügt
\end{itemize}

\textbf{Aufrufbeispiele aus \texttt{druckeErgebnis()}:}
\begin{lstlisting}[language=Java]
// Ankunft-Zeile Hinfahrt
printRow("An ", n -> n.getAnkunfthin() != null
    ? String.format("%02d", n.getAnkunfthin()) : "  ", plan, false);

// Stationszeile (mit Kollisions-x)
printRow("   ", n -> String.format("%2s", n.getName()), plan, true);

// Wartezeit-Zeile Rückfahrt
printRow("Wa ", n -> n.getWartezeitRueck() > 0
    ? String.format("%02d", n.getWartezeitRueck()) : "  ", plan, false);
\end{lstlisting}

% ────────────────────────────────────────────────────────────────────────────

\subsection{\texttt{EingabeLeser.leseDatei()}}

\begin{tcolorbox}[colback=gray!5, colframe=gray!40, title=Methodensignatur]
\texttt{public static SimulationsDaten leseDatei(String dateiPfad) throws IOException}
\end{tcolorbox}

\textbf{Zweck:} Liest eine Textdatei im vorgeschriebenen Format ein und baut
daraus ein \texttt{SimulationsDaten}-Objekt auf.

\textbf{Algorithmus (State-Machine über Zeilenindex):}

\begin{enumerate}
  \item Lese alle Zeilen mit \texttt{Files.readAllLines(Path.of(dateiPfad))}
  \item Iteriere mit Index $i$ über alle Zeilen:
    \begin{itemize}
      \item Zeile $i$ == \texttt{"Strecke:"} → Zeile $i+1$ mit \texttt{split("\\s+")}
            aufteilen; je Name \texttt{addBahnhof()} aufrufen
      \item Zeile $i$ == \texttt{"Abstaende:"} → Zeile $i+1$ in \texttt{int[]}
            umwandeln; \texttt{setAbstaende()} aufrufen
      \item Zeile $i$ == \texttt{"Start Hinfahrt:"} → Zeile $i+1$ als
            \texttt{Integer.parseInt()} lesen
    \end{itemize}
  \item Rückgabe: \texttt{new SimulationsDaten(bahnhofPlan, startZeit)}
\end{enumerate}

\textbf{Robustheit:} Die Methode verwendet \texttt{trim()} auf jede gelesene
Zeile, um führende/nachfolgende Leerzeichen zu entfernen. Unbekannte Zeilen
werden still ignoriert (kein Exception-Wurf bei leeren Zeilen).

\textbf{Einschränkung (BUG-02):} \texttt{startZeit} wird nicht auf $[0,59]$
normalisiert. Der Wert wird direkt von \texttt{Integer.parseInt()} übernommen.

% ─────────────────────────────────────────
\section{Zusammenfassung}

Das Programm setzt das \textbf{Strategy Pattern} ein, um drei verschiedene Fahrplanberechnungsalgorithmen hinter einer einheitlichen Schnittstelle (\texttt{FahrplanStrategie}) zu kapseln. Die Kernalgorithmen werden als gemeinsame statische Methoden in der Basisklasse implementiert und von den konkreten Strategien wiederverwendet.

Die Bahnhöfe werden als \textbf{doppelt verkettete Liste} modelliert. Diese Datenstruktur ermöglicht eine natürliche Traversierung in beide Richtungen (Hin- und Rückfahrt) und eine lokale Kollisionsprüfung, die nur den aktuellen Knoten und seinen direkten Nachbarn benötigt.

Ein wichtiges Entwurfsmerkmal ist die \textbf{gemeinsame Objektreferenz}: Alle drei Strategien operieren auf demselben \texttt{BahnhofPlan}-Objekt. Die \texttt{EinfacheFahrtStrategie} schreibt die Hinfahrtdaten in die Knoten, und die nachfolgenden Strategien lesen diese Daten direkt, ohne die Hinfahrt neu zu berechnen. Dies funktioniert korrekt, weil Strategien 2 und 3 die Hinfahrtdaten unverändert lassen.

Die \textbf{Kollisionserkennung} basiert auf einer Intervall-Überschneidungsprüfung auf einer Modulo-60-Uhr, die auch den Übergang von Minute 59 auf Minute 0 korrekt behandelt.

Die \textbf{Ausgabe} nutzt ein Funktionsinterface (\texttt{Function<BahnhofNode, String>}), um mit einer einzigen \texttt{printRow()}-Hilfsmethode alle sieben Zeilen der tabellarischen Darstellung zu erzeugen, ohne Code zu duplizieren.

\end{document}
